\chapter{数值表示与算术电路}

\begin{keyidea}
硬件只搬运固定宽度 bit，数值含义来自解释规则。本章从补码和标志位进入加法、移位、比较、乘法与 ALU 的实现取舍。
\end{keyidea}

\begin{classicnote}
经典体系结构教材讲 ALU 和数据通路时，很少把它们当作孤立名词处理，而是反复追问三件事：输入 bit 的解释是什么，组合路径的输出是什么，状态在什么边界提交。若能为一条动作写出“源寄存器、选择器、ALU 操作、标志位、写回来源、写使能、下一状态”的 trace（执行轨迹），就已经具备走读最小 CPU 的基本能力。
\end{classicnote}

\begin{pdfdiagram}
  \node[hw state, minimum width=2.3cm] (bits) at (0,0) {固定宽度 bit};
  \node[hw compute, minimum width=2.1cm] (alu) at (2.9,0) {ALU};
  \node[hw compute, minimum width=2.15cm] (flags) at (5.7,0) {标志位};
  \node[hw state, minimum width=2.35cm] (reg) at (8.65,0) {写回寄存器};
  \node[hw control, minimum width=2.1cm] (ctrl) at (4.35,-1.45) {控制器};
  \draw[hw bus] (bits) -- (alu);
  \draw[hw bus] (alu) -- (flags);
  \draw[hw bus] (flags) -- (reg);
  \node[hw label] at (1.45,0.58) {解释规则};
  \node[hw label] at (4.3,0.58) {结果与标志};
  \node[hw label] at (7.15,0.58) {条件判断};
  \draw[hw bus] (ctrl.north) -- (alu.south);
  \draw[hw bus] (ctrl.north east) -- (flags.south);
  \draw[hw bus] (ctrl.east) -| (reg.south);
  \node[hw label, text width=9.2cm] at (4.35,-2.35) {bit 解释、组合计算、标志生成和时钟提交连成同一条连贯路径。};
\end{pdfdiagram}
\diagramnote{ALU 章节的总图。箭头表示信息流向：固定 bit 先被解释，再进入组合计算，最后由控制器决定是否提交。}

\section{一、固定宽度 bit 的数值解释}

硬件只保存 bit。所谓整数、正负号、小数点、十进制数字、溢出、大小比较，都是对同一组 bit 模式的解释规则。若不先讲清楚表示，后面的加法器、减法器、比较器和 ALU 就会像在做抽象数学；实际上它们处理的始终是固定宽度的一组电线。位宽一旦确定，硬件能保存的模式数量也就确定：4-bit 有 16 种模式，8-bit 有 256 种模式，32-bit 有 $2^{32}$ 种模式。显示成十进制、十六进制还是二进制只是记法，不能改变硬件位宽。

4-bit 无符号数 \code{b3 b2 b1 b0} 的值是：

\begin{lstlisting}
value = b3*8 + b2*4 + b1*2 + b0
\end{lstlisting}

因此 \code{1011} 表示 11，因为它等于 8+0+2+1。无符号数没有负数概念，所有 bit 都贡献非负权重。4-bit 能表达的范围是 0 到 15，8-bit 能表达的范围是 0 到 255。一个 4-bit 加法器只能输出 4 个结果 bit，再加上可能的最高位外进位。若 15 加 1，低 4 位回到 \code{0000}，最高位外产生进位。这个进位不是“可有可无的第 5 位”，而是无符号运算超出 4-bit 范围时硬件给出的明确信号。

\begin{longtable}{L{0.18\linewidth}L{0.2\linewidth}L{0.28\linewidth}L{0.24\linewidth}}
\toprule
位宽 & 模式数量 & 无符号范围 & 典型硬件含义 \\
\midrule
4-bit & 16 & 0 到 15 & 一片 4-bit 加法器、一个十六进制数字 \\
8-bit & 256 & 0 到 255 & 早期微处理器常见数据宽度、一个字节 \\
16-bit & 65536 & 0 到 65535 & 8086 通用寄存器和 ALU 的核心宽度 \\
32-bit & $2^{32}$ & 0 到 $2^{32}-1$ & 80386 扩展后的通用寄存器与地址计算宽度 \\
\bottomrule
\end{longtable}

补码表示让负数也能走同一个加法器。以 4-bit 为例，最高位不再只是 8，而是解释成负权重 -8：

\begin{lstlisting}
value = b3*(-8) + b2*4 + b1*2 + b0
\end{lstlisting}

于是 \code{1111} 的值是 -8+4+2+1=-1，\code{1110} 的值是 -2，\code{1000} 的值是 -8。4-bit 补码范围是 -8 到 7。补码范围不对称，因为 \code{0000} 已经表示 0，负数一侧多出一个最小值。这个细节很重要：在固定宽度内，最小负数没有对应的正数。例如 4-bit 的 -8 取相反数仍会得到 \code{1000}，这不是数学规则失效，而是硬件位宽不足。

为什么负权重成立？补码不是人为规定“最高位算负数”，而是模运算的自然结果。4-bit 加法器算的是模 16 加法：结果只保留低 4 位，超出 16 的部分被丢弃。在模 16 的世界里，-3 和 13 是同一个数，因为 $13+3=16\equiv 0$。所以“存 -3”就是“存 13”，即 \code{1101}。同一根线上的 bit 模式于是有了两种读法：无符号读作 13，有符号读作 -3。把最高位解释成 -8 而不是 +8，两种读法就对上了：$8+4+1=13$，$-8+4+1=-3$，而 $13\equiv-3 \pmod{16}$。负权重不是新算术，只是给同一个模运算换了一个名字区间：0 到 15 换成 -8 到 7。

\begin{pdfdiagram}
  \draw[BookRule, line width=0.6pt] (0,0) circle (1.55);
  \foreach \i/\u/\s in {0/0/0, 1/1/1, 2/2/2, 3/3/3, 4/4/4, 5/5/5, 6/6/6, 7/7/7, 8/8/{-8}, 9/9/{-7}, 10/10/{-6}, 11/11/{-5}, 12/12/{-4}, 13/13/{-3}, 14/14/{-2}, 15/15/{-1}} {
    \draw[BookMuted, line width=0.5pt] ({90-22.5*\i}:1.42) -- ({90-22.5*\i}:1.68);
    \node[font=\scriptsize\sffamily, text=BookMuted] at ({90-22.5*\i}:1.98) {\u};
    \node[font=\scriptsize\sffamily\bfseries, text=BookInk] at ({90-22.5*\i}:1.06) {\s};
  }
  \draw[fill=BookRed] ({90-22.5*13}:1.55) circle (0.09);
  \node[hw label, text=BookRed, text width=4.6cm, align=left] at (4.0,-1.1) {\code{1101}：无符号读 13\\补码读 -3\\$13+3=16\equiv0$};
  \node[hw label, align=center] at (0,-2.55) {4-bit 模 16 数环：外圈无符号，内圈补码};
\end{pdfdiagram}
\diagramnote{补码的数环读法：同一个 bit 模式，外圈是无符号读数，内圈是补码读数。加法器只按模 16 进位，两种读法共享同一条进位链——这就是为什么加法和减法能共用一套硬件。}

“取反加一”有一行证明：对任意 $x$，都有 $x + \sim x = \code{1111} = -1$（模 16 下各位互补，相加必得全 1），所以 $-x = \sim x + 1$。硬件只需要一排 NOT 门和一条加法进位链，减法就被归并成了加法。

求一个补码数的相反数，可以按“取反加一”理解：

\begin{lstlisting}
x      = 0011  // 3
~x     = 1100
~x + 1 = 1101  // -3 in 4-bit two's complement
\end{lstlisting}

这个规则适合硬件，因为取反可以用一排 NOT 或 XOR 门完成，加一可以交给加法器。于是减法可以变成加法：

\begin{lstlisting}
a - b = a + (NOT b) + 1
\end{lstlisting}

这条公式是 ALU 设计的关键。硬件不必为加法和减法各放一套完全独立的算术电路，只要让控制信号同时决定 B 操作数是否反相，以及最低位进位输入是否为 1，就能在同一条进位链上完成 ADD 和 SUB：\hwkey{控制信号改变的是电路连接和输入选择，不是纸面公式}。

\topic{原码、反码、补码是三种被真实采用过的负数表示}
补码不是唯一被试过的方案。早期计算机还认真使用过两种更直觉的表示：原码把最高位专门用作符号位，其余位保存绝对值；反码把正数逐位取反得到对应负数；补码则在反码基础上再加一。三种方案用 8-bit 都能表示 -5，但硬件代价差别很大。先把同一个负数的三种写法并排摆开：

\begin{longtable}{L{0.12\linewidth}L{0.28\linewidth}L{0.16\linewidth}L{0.32\linewidth}}
\toprule
表示法 & 负数规则 & -5 的写法 & 零的写法 \\
\midrule
原码 & 最高位记符号，其余位记绝对值 & \code{1000 0101} & \code{0000 0000} 与 \code{1000 0000} 两个零 \\
反码 & 正数逐位取反 & \code{1111 1010} & \code{0000 0000} 与 \code{1111 1111} 两个零 \\
补码 & 逐位取反再加一 & \code{1111 1011} & 唯一的 \code{0000 0000} \\
\bottomrule
\end{longtable}

补码最终胜出，三条原因都直接对应门电路的开销。第一，加法会自然处理符号。原码做加减前必须先比较两数符号和绝对值大小，决定实际执行加法还是减法，再单独确定结果符号——符号位不能参与进位，需要一条独立的逻辑路径。补码把符号位当成普通位一起进位：\code{1111 1011 + 0000 0101 = 1 0000 0000}，最高位进位被位宽丢弃，低 8 位就是 0，(-5)+5 在普通加法器上一次算对，不需要任何符号特例。第二，零唯一。原码和反码都有 +0 和 -0 两种模式，“结果是否为零”要检测两种 bit 模式，判零电路多一级组合逻辑；补码只有一种零，Zero 标志就是对所有结果 bit 做一次 NOR。第三，取负便宜且单向。反码取负确实只要一排 NOT 门，但它的加法在产生最高位进位时必须把进位绕回最低位再加一次，这叫循环进位（end-around carry），进位链因此多绕一圈；补码取负是“取反加一”，NOT 加已有进位链即可完成，进位始终朝一个方向走。

循环进位的差别用 (-5)+7 演算一次。反码：\code{1111 1010 + 0000 0111 = 1 0000 0001}，最高位的进位 1 必须加回最低位，\code{0000 0001 + 1 = 0000 0010}，才得到 2。补码：\code{1111 1011 + 0000 0111 = 1 0000 0010}，进位直接丢弃，低 8 位 \code{0000 0010} 就是 2。一次加法里反码要等进位回绕，补码只走一趟。

这三条性质的意义超出表示法本身：补码把符号处理、判零和取负全部归并到同一条进位链上，加减法、比较和取负共用一套硬件。历史上 UNIVAC 1100 系列等机器沿用反码，而 x86、ARM、RISC-V、MIPS 的有符号整数一律是补码。前面“减法复用加法器”的全部便利，本质上都来自这张表最后一行。

同一组 bit 可以按无符号解释，也可以按补码解释。因此溢出判断也不同。

\begin{longtable}{L{0.2\linewidth}L{0.34\linewidth}L{0.34\linewidth}}
\toprule
解释方式 & 关注点 & 例子 \\
\midrule
无符号 & 是否产生最高位外的进位 & 15 + 1 超出 4-bit 无符号范围 \\
补码有符号 & 正正得负或负负得正 & 7 + 1 得到 \code{1000}，被解释为 -8 \\
\bottomrule
\end{longtable}

以下 4-bit 加法给出同一结果在两种解释下的差异：

\begin{lstlisting}
0111 + 0001 = 1000
\end{lstlisting}

若按无符号解释，这是 7+1=8，仍在表示范围内。若按补码解释，\code{0111} 是 7，\code{0001} 是 1，结果 \code{1000} 是 -8，正数加正数却得到负数，所以发生有符号溢出。硬件不会替使用者决定哪种解释成立；它只提供结果 bit、最高位外进位、符号位和溢出标志。程序、指令语义和控制器必须约定当前采用哪一种解释。

\topic{溢出判定演算：四个 8-bit 例子覆盖全部标志组合}
把“正正得负、负负得正”推进到 8-bit，用四个例子走完整数值过程。8-bit 补码范围是 -128 到 127，无符号范围是 0 到 255：

\begin{lstlisting}
+100 +  +50: 0110 0100 + 0011 0010 =  1001 0110
             signed: +150 read as -106   WRONG (overflow)
             Cout = 0, Overflow = 1

-100 +  -50: 1001 1100 + 1100 1110 =1 0110 1010
             signed: -150 read as +106   WRONG (overflow)
             Cout = 1, Overflow = 1

+100 +  -50: 0110 0100 + 1100 1110 =1 0011 0010
             signed: +50                 correct
             Cout = 1, Overflow = 0

-100 +  +50: 1001 1100 + 0011 0010 =  1100 1110
             signed: -50                 correct
             Cout = 0, Overflow = 0
\end{lstlisting}

判定规则可以总结成一句：\hwkey{两个同号数相加，得到异号结果，即补码溢出}。写成布尔式，若 \code{a7}、\code{b7}、\code{r7} 是两个输入和结果的符号位，则 $V = a_7 b_7 \overline{r_7} + \overline{a_7}\,\overline{b_7}\, r_7$。硬件里还有一个更便宜的等价判法：进入符号位的进位 \code{C7} 与最高位向外的进位 \code{Cout} 不一致就是溢出，$V = C_7 \oplus C_{out}$。用第一个例子验证：逐位进位（\code{c1} 到 \code{c8}）是 \code{0,0,0,0,0,1,1,0}，进入符号位的进位是 1、向外进位是 0，两者不同，$V=1$，与符号判法一致。

为什么异号相加永不溢出？一个正数加一个负数，结果必然落在两个加数之间：和的绝对值不超过两数绝对值中较大的那个，而那个加数本身就在 -128 到 127 之内。范围内加范围内得到范围外，只有在两个加数同号、绝对值同向叠加时才可能发生。减法的判定也归并到同一条规则：把它看成加上取负后的数，于是“异号相减且结果符号与被减数不同”就是溢出。

Carry 和 Overflow 是两个相互独立的判断，上面四个例子恰好把 \code{(Cout, V)} 的四种组合 \code{(0,1)}、\code{(1,1)}、\code{(1,0)}、\code{(0,0)} 都走了出来。Carry 报告的是“bit 模式超出无符号范围”，即模 $2^{8}$ 回绕发生；Overflow 报告的是“补码真值超出有符号范围”。第二个例子里两者同时为 1：无符号读是 $156+206=362$ 超界，补码读是 $-100+(-50)=-150$ 超界。第三个例子里 Carry=1 但 V=0：无符号读 $100+206=306$ 回绕，补码读 $+50$ 完全正确。硬件不替程序选择哪种读法，所以两个标志都必须产生。

\begin{longtable}{L{0.16\linewidth}L{0.14\linewidth}L{0.13\linewidth}L{0.1\linewidth}L{0.08\linewidth}L{0.27\linewidth}}
\toprule
演算 & 输入符号 & 结果符号 & Cout & V & 读法结论 \\
\midrule
+100 + +50 & 同号（正） & 负 & 0 & 1 & 补码溢出；无符号 150 仍正确 \\
-100 + -50 & 同号（负） & 正 & 1 & 1 & 补码溢出；无符号 362 也超界 \\
+100 + -50 & 异号 & 正 & 1 & 0 & 补码正确；无符号发生回绕 \\
-100 + +50 & 异号 & 负 & 0 & 0 & 两种读法都正确 \\
\bottomrule
\end{longtable}

两个标志的用途也因此分开：多字长加法把低字节的 Carry 送进高字节的 \code{Cin}，有符号范围检查看 Overflow。x86-64 把两者分别放在 CF 和 OF，条件转移各走各的：

\begin{lstlisting}
; x86-64, Intel syntax
mov al, 100
add al, 50     ; AL = 0x96, CF = 0, OF = 1
jo  too_big    ; signed path checks OF

mov al, 200
add al, 100    ; AL = 0x2C, CF = 1, OF = 0
jc  carry_in   ; unsigned path checks CF
\end{lstlisting}

这段分工的意义在于：标志位不是“结果好或坏”的笼统指示，而是两种解释规则各自的边界信号。控制器和编译器按当前约定取其中一个，另一个照常产生但可以被忽略。

位宽变化也要有明确规则。把 4-bit 值送入 8-bit 数据通路时，若它是无符号数，应做零扩展；若它是补码有符号数，应做符号扩展。把 8-bit 值截断成 4-bit 时，高位被丢弃，剩余低 4 位不一定仍表示原值。

\begin{longtable}{L{0.16\linewidth}L{0.27\linewidth}L{0.27\linewidth}L{0.2\linewidth}}
\toprule
动作 & 输入解释 & 输出规则 & 示例 \\
\midrule
零扩展 & 无符号 & 高位补 0 & \code{1011 -> 00001011} \\
符号扩展 & 补码有符号 & 高位复制符号位 & \code{1011 -> 11111011} \\
截断 & 固定低位 & 只保留低位 & \code{10011011 -> 1011} \\
饱和 & 特定 DSP（数字信号处理器）/多媒体语义 & 超出范围钳到最大/最小 & 需要额外比较和选择 \\
\bottomrule
\end{longtable}

符号扩展不是显示格式，而是真实硬件路径。若 4-bit 的 \code{1011} 按补码解释为 -5，扩展到 8-bit 后必须成为 \code{11111011}，仍然表示 -5；若错误地零扩展为 \code{00001011}，它就变成 11。这个错误在地址偏移、短立即数、分支位移和加载有符号字节时都可能出现。8086、80386 这类处理器的指令语义中会明确区分零扩展、符号扩展和普通截断，因为这些规则直接决定 ALU 输入。

\topic{字节、字和字节序把数值解释放进内存}
体系结构教材常从“一个对象在内存中占哪些字节”开始训练机器级思维。硬件寄存器里可以一次保存 16-bit、32-bit 或 64-bit 值，但主存通常按字节编址。于是一个 32-bit 数值放进内存时，必须规定四个字节按什么顺序排列。若数值为 \code{0x12345678}，大端序把最高有效字节放在最低地址，小端序把最低有效字节放在最低地址。两者保存的是同一个 32-bit 模式，只是地址到字节车道的映射不同。

\begin{longtable}{L{0.18\linewidth}L{0.32\linewidth}L{0.4\linewidth}}
\toprule
字节序 & 低地址到高地址 & 硬件含义 \\
\midrule
大端序 & \code{12 34 56 78} & 低地址保存最高有效字节，适合按书写顺序观察十六进制 \\
小端序 & \code{78 56 34 12} & 低地址保存最低有效字节，低位字节扩展和多精度运算方便 \\
\bottomrule
\end{longtable}

字节序不是“显示问题”，而是 LOAD/STORE、字节使能和总线车道的硬件约定。CPU 执行 32-bit LOAD 时，内存系统返回四个相邻字节，处理器内部再按本机字节序把它们拼成寄存器值。CPU 执行 8-bit LOAD 时，只取其中一个字节；若目标寄存器更宽，还要按指令语义做零扩展或符号扩展。若把这三件事混在一起，调试时会出现非常典型的错误：内存窗口看见的字节顺序和寄存器中显示的十六进制值不一致，看起来像“读错了”，实际可能只是字节序解释不同。

\begin{lstlisting}
memory bytes at increasing addresses:
  0x1000: 78
  0x1001: 56
  0x1002: 34
  0x1003: 12

little-endian 32-bit load from 0x1000 -> 0x12345678
8-bit load from 0x1000                -> 0x78
\end{lstlisting}

不同宽度的访问有各自必须写清的规则。字节 LOAD 选择一个字节车道，再扩展到寄存器宽度——零扩展还是符号扩展必须由指令规定；半字 LOAD 选择两个相邻字节并按字节序拼接，地址低位和对齐规则必须满足；字 LOAD 选择四个字节拼成 32-bit 值，小端/大端、字节使能和未对齐处理都要写清；字节 STORE 只打开一个字节写使能，绝不能破坏同一字中的其他字节。

\topic{内存对象和指针本质上也是位模式}
从硬件角度看，指针不是神秘类型，而是一个被解释为地址的固定宽度 bit 模式。数组、结构体、栈帧和缓冲区最终都归结为“起始地址 + 偏移”的地址形成路径。若数组元素是 4 字节整数，访问 \code{a[i]} 时地址形成部件通常要计算 \code{base + i*4}；若结构体字段在对象内偏移 12 字节，LOAD/STORE 的地址就是 \code{base + 12}。这与前面的定点数类似：bit 本身不变，解释约定决定它是整数、地址、偏移还是控制字段。

把几类对象的硬件路径分开看：标量整数是固定位宽 bit 模式，存在寄存器或内存字节里，经 ALU 解释；指针是地址宽度 bit 模式，会进入地址形成、译码、cache/TLB 或总线；数组元素的地址是基址加缩放下标，由加法器、移位器或地址生成单元算出；结构体字段的地址是基址加固定偏移，通常由立即数扩展和 ALU 地址加法完成；栈对象的地址是栈指针加偏移，依赖 SP/RSP 类寄存器和调用约定。

这个视角能解释为什么越界访问危险。硬件不会自动知道“这个地址还属于数组”。若没有保护机制，\code{a[10]} 只是另一次地址加法和一次 LOAD/STORE，它可能指向相邻变量、返回地址、MMIO（内存映射 I/O）窗口或未映射区域。分页、IOMMU 和异常机制本质上就是把“哪些地址可访问、按什么权限访问”提升为硬件可检查的约定。

\topic{小数不是新电路，而是新的位权解释}整数位的权重是 $1,2,4,8,\ldots$，小数位只是把二进制小数点右侧的权重继续向下分成 $1/2,1/4,1/8,\ldots$。例如：

\begin{lstlisting}
10.101_2 = 1*2 + 0*1 + 1*(1/2)
         + 0*(1/4) + 1*(1/8)
         = 2.625
\end{lstlisting}

这里的“小数点”不是一根额外电线。若寄存器里只保存 \code{10101}，硬件本身看见的仍然是五个 bit；把它解释成 $21$、$10.101_2$，还是某种编码字段，取决于电路和指令约定。这个事实会直接影响硬件设计：二进制小数的加法仍可以复用整数加法器，但必须保证两个操作数的小数点位置相同；若小数点位置不同，必须先对齐。

二进制小数和十进制小数并不总能互相有限表示。$0.5$ 可以写成 $0.1_2$，因为它正好是 $1/2$；$0.25$ 可以写成 $0.01_2$，因为它正好是 $1/4$。十进制的 $0.1$ 则不能用有限个二进制小数位精确表示，只能形成无限循环展开。计算器、编译器和 FPU（浮点运算单元）都受位宽限制：寄存器只有固定数量的 bit，必须在某个位置舍入。

\begin{longtable}{L{0.16\linewidth}L{0.24\linewidth}L{0.24\linewidth}L{0.26\linewidth}}
\toprule
十进制值 & 二进制表示 & 是否有限 & 硬件含义 \\
\midrule
0.5 & $0.1_2$ & 有限 & 一位小数即可精确保存 \\
0.25 & $0.01_2$ & 有限 & 权重为 $1/4$ 的 bit 置 1 \\
0.625 & $0.101_2$ & 有限 & $1/2+1/8$ \\
0.1 & 无限循环 & 不能有限表示 & 固定位宽内必须近似和舍入 \\
\bottomrule
\end{longtable}

\topic{定点数把小数点位置变成硬件约定}定点数把“小数点在第几位”固定下来。常见写法是 Q 格式（沿用早期 DSP 手册用字母 Q 标记小数位数的习惯）：Q4.4 表示 8-bit 总宽度中有 4 个整数相关位和 4 个小数位，实际值等于原始整数除以 $2^4$。例如 \code{00011000} 的无符号整数值是 24；若按 Q4.4 解释，它的实际值是 $24/16=1.5$。同一串 bit 并没有改变，改变的是缩放约定。

把 \code{00011000} 的每一位按 Q4.4 的权重新画一遍，缩放约定就看得更具体：

\begin{pdfdiagram}
  % 位权标注（上排）：Q4.4 约定下各 bit 的权重
  \node[hw label] at (0.55,1.5) {$2^{3}$};
  \node[hw label] at (1.65,1.5) {$2^{2}$};
  \node[hw label] at (2.75,1.5) {$2^{1}$};
  \node[hw label] at (3.85,1.5) {$2^{0}$};
  \node[hw label] at (4.95,1.5) {$2^{-1}$};
  \node[hw label] at (6.05,1.5) {$2^{-2}$};
  \node[hw label] at (7.15,1.5) {$2^{-3}$};
  \node[hw label] at (8.25,1.5) {$2^{-4}$};
  % 8 个 bit 格：0001 1000
  \draw[BookTeal, line width=0.9pt] (0,0.5) rectangle (1.1,1.2);
  \node at (0.55,0.85) {0};
  \draw[BookTeal, line width=0.9pt] (1.1,0.5) rectangle (2.2,1.2);
  \node at (1.65,0.85) {0};
  \draw[BookTeal, line width=0.9pt] (2.2,0.5) rectangle (3.3,1.2);
  \node at (2.75,0.85) {0};
  \draw[BookTeal, line width=0.9pt] (3.3,0.5) rectangle (4.4,1.2);
  \node at (3.85,0.85) {1};
  \draw[BookTeal, line width=0.9pt] (4.4,0.5) rectangle (5.5,1.2);
  \node at (4.95,0.85) {1};
  \draw[BookTeal, line width=0.9pt] (5.5,0.5) rectangle (6.6,1.2);
  \node at (6.05,0.85) {0};
  \draw[BookTeal, line width=0.9pt] (6.6,0.5) rectangle (7.7,1.2);
  \node at (7.15,0.85) {0};
  \draw[BookTeal, line width=0.9pt] (7.7,0.5) rectangle (8.8,1.2);
  \node at (8.25,0.85) {0};
  % bit 索引（下排）
  \node[hw label] at (0.55,0.2) {bit7};
  \node[hw label] at (1.65,0.2) {bit6};
  \node[hw label] at (2.75,0.2) {bit5};
  \node[hw label] at (3.85,0.2) {bit4};
  \node[hw label] at (4.95,0.2) {bit3};
  \node[hw label] at (6.05,0.2) {bit2};
  \node[hw label] at (7.15,0.2) {bit1};
  \node[hw label] at (8.25,0.2) {bit0};
  % 小数点约定位置
  \draw[hw return] (4.4,-0.25) -- (4.4,1.95);
  \node[hw label, anchor=west] at (4.55,1.95) {小数点约定位置（不占 bit）};
  % 两种读法
  \node[hw label, anchor=west] at (0,-0.75) {无符号整数读法：bit4、bit3 置 1，$16+8=24$};
  \node[hw label, anchor=west] at (0,-1.3) {Q4.4 读法：$2^{0}+2^{-1}=1+0.5=1.5$，即 $24/2^{4}=24/16$};
\end{pdfdiagram}
\diagramnote{同一个 8-bit 模式 \code{00011000} 在 Q4.4 约定下的位权解读。电路里的 8 个 bit 不变，约定把小数点固定在 bit 4 与 bit 3 之间：高 4 位权重 $2^{3}$ 到 $2^{0}$，低 4 位权重 $2^{-1}$ 到 $2^{-4}$。置 1 的是 bit4 与 bit3，按整数读为 $16+8=24$，按 Q4.4 读为 $1+0.5=1.5$——两种读法相差的正是缩放因子 $2^{4}$，即 $24/16=1.5$。}

\begin{longtable}{L{0.18\linewidth}L{0.18\linewidth}L{0.24\linewidth}L{0.3\linewidth}}
\toprule
原始 bit & 原始整数 & 定点解释 & 说明 \\
\midrule
\code{00011000} & 24 & Q4.4 中为 1.5 & 小数位数为 4，除以 16 \\
\code{00000110} & 6 & Q1.7 中为 0.046875 & 小数位数为 7，除以 128 \\
\code{11110000} & -16 & 有符号 Q4.4 中为 -1.0 & 先按补码得 -16，再除以 16 \\
\code{01111111} & 127 & 有符号 Q1.7 中为 0.9921875 & 最大正值小于 1 \\
\bottomrule
\end{longtable}

定点加减最适合用现有整数 ALU。只要两个数采用相同 Q 格式，硬件可以直接把原始 bit 送入加法器；结果仍按同一个缩放因子解释。例如 Q4.4 中的 $1.5+0.25$ 可以看成原始整数 $24+4=28$，结果 \code{00011100} 解释为 $28/16=1.75$。因此许多音频、控制、电机、传感器和早期图形场景会优先使用定点数：整数加法器、移位器和乘法器即可承担主要工作。

定点乘法需要多一步处理。若两个 Q4.4 数相乘，原始整数相乘后缩放因子变成 $2^8$，因为两个 $2^4$ 相乘得到 $2^8$。要把结果放回 Q4.4，通常要右移 4 位，并决定被丢弃的低位是截断、四舍五入还是参与饱和判断。

\begin{lstlisting}
Q4.4 value 1.5  -> raw 24
Q4.4 value 0.5  -> raw 8
raw product     -> 24 * 8 = 192
rescale to Q4.4 -> 192 >> 4 = 12
Q4.4 result     -> 12 / 16 = 0.75
\end{lstlisting}

这个例子说明定点是一种明确的硬件选择。它牺牲动态范围，换来较低延迟（时延）、较小面积和更容易预测的执行时间。若电路要实时控制电机电流、滤波音频采样或处理传感器数据，固定缩放因子往往比复杂浮点更合适。

\begin{chipcase}
许多微控制器并不一定带硬件 FPU，或者即使带 FPU，也会在中断、功耗和实时性严格的路径上使用整数和定点运算。DSP 芯片则常见乘加单元，也就是 MAC 路径，用于反复执行“乘法、对齐、累加、饱和”。这些芯片的手册通常会明确给出 Q 格式、饱和模式、舍入模式和累加器宽度，因为它们决定滤波器、控制环和音频处理是否稳定。
\end{chipcase}

\topic{定点乘法先做整数乘，再把小数点移回约定位置}
两个 Q4.4 数相乘时，硬件看见的仍然只是两个 8-bit 整数。完整走一遍 $1.5 \times 2.25$：

\begin{lstlisting}
1.5   -> raw 24 = 0001 1000     (24 / 16 = 1.5)
2.25  -> raw 36 = 0010 0100     (36 / 16 = 2.25)

integer product: 24 * 36 = 864
16-bit accumulator: 0000 0011 0110 0000
read as Q8.8: 864 / 256 = 3.375

rescale to Q4.4: 864 >> 4 = 54 = 0011 0110
check: 54 / 16 = 3.375 = 1.5 * 2.25
\end{lstlisting}

为什么乘完必须移动小数点？每个操作数的 raw 值都是“真值 $\times 2^{4}$”，整数乘积天然携带 $2^{4} \times 2^{4} = 2^{8}$ 的缩放。若把 864 直接按 Q4.4 读，得到的是 54，恰好是真值的 16 倍。小数点在电路里不是实体，而是“缩放因子是多少”的约定；乘法把两个约定相乘，所以必须右移 4 位把约定改回 Q4.4。若目标是更宽的累加器，也可以不立即移位，按 Q8.8 继续参与后续乘加，等写回存储前再统一缩放——缩放时机是设计选择，缩放动作不可省略。

被右移丢弃的低位是真实的信息损失，处理方式必须写清。直接 \code{>>4} 是截断：$1.1875 \times 1.1875$ 的 raw 乘积是 $19 \times 19 = 361$，\code{361 >> 4 = 22}，即 1.375，而真值是 1.41015625。若要在 Q4.4 内就近舍入，右移前先加半个最低位的权重，即加 \code{1000_2}：\code{(361 + 8) >> 4 = 23}，即 1.4375，距真值更近。两种做法不是对错问题，而是误差分布的约定：控制环在意单向误差累积，音频在意舍入噪声，芯片手册会把所选模式写明。

溢出检查同样在移位这一步完成。两个 8-bit 数相乘必须用 16-bit 暂存；右移 4 位后，若只取低 8 位作为 Q4.4 结果，则 bit 15 到 bit 8 必须全是 bit 7（新符号位）的复制，否则值已经超出 8-bit Q4.4 的范围。例如 $7.9375 \times 7.9375$：raw 乘积 $127 \times 127 = 16129$ = \code{0011 1111 0000 0001}，右移 4 位得 \code{0000 0011 1111 0000}；此时 bit 7 为 1，而 bit 15 到 bit 8 是 \code{0000 0011}，不是全 1，判定溢出——真值约 63.004 远超有符号 Q4.4 的上限 7.9375。此时要么置溢出标志交给软件，要么按饱和约定钳到 \code{0111 1111}（7.9375）。

\begin{longtable}{L{0.16\linewidth}L{0.3\linewidth}L{0.44\linewidth}}
\toprule
步骤 & 数值或位模式 & 缩放约定与检查点 \\
\midrule
解释输入 & raw 24 与 raw 36 & 各带 $2^{4}$ 缩放，两数小数位数必须一致 \\
整数相乘 & $24 \times 36 = 864$ & 缩放变为 $2^{8}$，需要双倍位宽暂存 \\
右移还原 & \code{864 >> 4 = 54} & 回到 Q4.4；右移前加 8 可就近舍入 \\
溢出检查 & 高位必须只是符号复制 & 超界时置标志，或饱和到最大/最小值 \\
\bottomrule
\end{longtable}

真实 DSP 把这条路径做成了专门部件。MAC（乘加）单元内部保留比输出格式宽得多的累加器，例如 32-bit 或 40-bit，让一连串“乘、加、乘、加”在中间过程不右移、不丢位，只在最终写回时做一次缩放、舍入和溢出检查。累加器多出来的高位常叫保护位（guard bits），作用正是推迟溢出判定：中间和允许暂时超出输出范围，只要最终结果能装回目标格式。这再次说明定点不是“简化版浮点”，而是一套把缩放时机、舍入模式和溢出策略全部显式写出来的工程约定。

\topic{十进制与 BCD 解决的是输入输出和精确十进制表示}二进制适合硬件算术，但人们日常习惯使用十进制。BCD，即 binary-coded decimal，用 4 个 bit 保存一个十进制数字。最常见的 8421 BCD 中，\code{0000} 到 \code{1001} 表示 0 到 9，\code{1010} 到 \code{1111} 不是合法十进制数字。两位十进制 59 可以保存为 \code{0101 1001}，而不是普通二进制整数 59 的 \code{00111011}。

\begin{longtable}{L{0.18\linewidth}L{0.28\linewidth}L{0.44\linewidth}}
\toprule
表示 & 示例 & 适用问题 \\
\midrule
普通二进制 & 59 为 \code{00111011} & ALU 运算紧凑，位宽利用率高 \\
8421 BCD & 59 为 \code{0101 1001} & 每个十进制位可直接显示和校正 \\
非法 BCD & \code{1010} 到 \code{1111} & 需要检测，不能当作十进制数字 \\
\bottomrule
\end{longtable}

BCD 加法的关键在于十进制校正。若一个半字节相加结果超过 9，或者低半字节产生了十进制意义上的进位，需要加 \code{0110} 把结果调整回合法 BCD。例如 $8+7=15$，普通二进制低 4 位是 \code{1111}，这不是合法 BCD 数字；加 \code{0110} 后得到 \code{1 0101}，表示十进制进位 1 和低位数字 5。早期处理器保留辅助进位、十进制调整等机制，往往就是为了兼容十进制输入输出、计算器和商业数据处理。

用 $59+38$ 走一遍完整流程，两个半字节的判断结果正好相反：

\begin{pdfdiagram}
  \node[hw state] (ops) at (0,0) {59 = 0101 1001，38 = 0011 1000（8421 BCD）};
  \node[hw compute] (add) at (0,-1.5) {低半字节二进制加：1001 + 1000 = 1 0001};
  \node[hw control] (dec) at (0,-3.0) {低半字节和超 9？本例 17 > 9：是};
  \node[hw compute] (fix) at (0,-4.5) {加 0110 校正：0001 + 0110 = 0111，十进制进位 1};
  \node[hw compute] (hi) at (0,-6.0) {高半字节加进位：0101 + 0011 + 1 = 1001 $\leq$ 9，不校正};
  \node[hw state] (res) at (0,-7.5) {结果 1001 0111 = 十进制 97};
  \draw[hw bus] (ops.south) -- (add.north);
  \draw[hw bus] (add.south) -- (dec.north);
  \draw[hw bus] (dec.south) -- node[hw label, right=2pt, pos=0.5]{是} (fix.north);
  \draw[hw bus] (fix.south) -- (hi.north);
  \draw[hw bus] (hi.south) -- (res.north);
  % 否分支：不超 9 则跳过校正，从右侧通道直接进入高半字节
  \draw[hw bus] (dec.east) -- node[hw label, above=2pt, pos=0.5]{否：不校正} (5.6,-3.0) -- (5.6,-6.0) -- (hi.east);
\end{pdfdiagram}
\diagramnote{BCD 加法的十进制校正流程，以 $59+38$ 为例。两个 8421 BCD 数先按普通二进制逐半字节相加；低半字节的和超 9（或产生半字节进位）时加 \code{0110}，把结果推回合法 BCD 并向上送出十进制进位，否则原样保留。本例低半字节 $9+8=17$ 需要校正，得数字 7 与进位 1；高半字节 $5+3+1=9$ 不再超 9，无需二次校正，最终 \code{1001 0111} 即十进制 97。}

\begin{chipcase}
8086 继承了面向 BCD 的十进制调整指令和辅助进位标志，原因不在于 BCD 比二进制 ALU 更快，而在于早期商业计算、显示和键盘输入大量使用十进制数字。现代通用 CPU 的主路径通常不会为 BCD 牺牲整数流水线，但十进制浮点、金融数据库和显示驱动仍会在特定边界上关心十进制精确性。
\end{chipcase}

\topic{浮点数把范围和精度分开}定点数的小数点位置固定，动态范围有限。浮点数采用另一种约定：用一个符号位表示正负，用指数字段表示尺度，用尾数（significand，也称有效数字）字段表示精度。可以把它理解成硬件版科学计数法。一个常见规格化二进制浮点数可写成：

\begin{lstlisting}
value = (-1)^sign * 1.fraction * 2^(exponent - bias)
\end{lstlisting}

其中 \code{bias} 是阶码偏置，用来让指数字段以无符号 bit 保存正负指数。尾数前面的 \code{1} 通常不显式保存，这叫隐藏位。浮点格式还要为特殊值保留编码：正负零、非规格化数（subnormal，也称次正规数）、正负无穷大和 NaN。非规格化数用于靠近 0 的渐进下溢；无穷大用于表示超出范围或除以 0 一类结果；NaN 表示“不是一个数”的异常传播，例如 $0/0$ 或无效运算。

把这条公式落实到最常见的格式上。IEEE 754 单精度用 32 个 bit，分成三段：

\begin{pdfdiagram}
  \draw[line width=0.55pt] (0,0) rectangle (12.4,0.9);
  \draw[line width=0.55pt] (0.8,0) -- (0.8,0.9);
  \draw[line width=0.55pt] (5.2,0) -- (5.2,0.9);
  \node at (0.4,0.45) {s};
  \node at (3.0,0.45) {阶码 e（8 位，bit 30–23）};
  \node at (8.8,0.45) {尾数 f（23 位，bit 22–0）};
  \node[hw label] at (3.0,-0.45) {按无符号保存，bias=127：e=127 表示 $2^{0}$};
  \node[hw label] at (8.8,-0.45) {规格化时前导 1 隐藏，不占 bit};
  \node[hw label] at (6.2,1.35) {值 = $(-1)^{s} \times 1.f \times 2^{e-127}$};
\end{pdfdiagram}
\diagramnote{IEEE 754 单精度的 32-bit 布局：1 位符号、8 位阶码、23 位尾数。阶码偏置 127，所以 e=1 表示 $2^{-126}$，e=254 表示 $2^{127}$；e=0 和 e=255 被保留给非规格化数、零、无穷大和 NaN。}

两个完整的解码演算，胜过十遍公式：

\begin{lstlisting}
0x3F800000:
  s = 0
  e = 0111 1111 = 127  ->  exponent = 127 - 127 = 0
  f = 000...0          ->  significand = 1.0
  value = +1.0 * 2^0 = 1.0

0x40490FDB:
  s = 0
  e = 1000 0000 = 128  ->  exponent = 1
  f = 100 1001 0000 1111 1101 1011
       significand = 1.5707963...
  value = 1.5707963 * 2^1 = 3.1415926...
\end{lstlisting}

对阶加法也用一个真实位宽走一遍。计算 $2^{10} + 2^{-14}$：两个数指数相差 24，较小数的尾数必须右移 24 位才能对齐；但单精度尾数只有 23 位，被移出的最低那个 1 落在 23 位尾数字段之外，由附加位（GRS）接住，按舍入规则处理——结果舍入回 $2^{10}$。$2^{-14}$ 这次加法在硬件里真实发生了，却被格式位宽“吃掉”了。

\begin{lstlisting}
2^10 + 2^-14 in binary32:
  align: shift small significand right by 24
  23-bit field cannot keep the shifted-out bit
  result rounds back to 2^10
\end{lstlisting}

四个字段各自的硬件动作也要分清：符号位进入符号处理和比较路径，不能直接等同整数的最高位；指数字段决定数值尺度，加减前要对阶、乘除时参与指数加减，风险是范围溢出或下溢；尾数是精度来源，进入尾数加法、乘法和规格化，风险是位数有限导致舍入；特殊值（零、非规格化、Inf、NaN）需要旁路和异常规则，比较、传播和转换时都要查清语义。

\topic{阶码的两个端点被保留给零、非规格化数、Inf 和 NaN}
单精度阶码 1 到 254 留给规格化数，全 0 和全 1 两个端点另有约定：\code{e = 0} 表示零和非规格化数，\code{e = 255} 表示无穷大和 NaN。于是一个 32-bit 模式先按阶码分五类，再看尾数细分：

\begin{longtable}{L{0.13\linewidth}L{0.12\linewidth}L{0.11\linewidth}L{0.32\linewidth}L{0.22\linewidth}}
\toprule
类别 & 阶码 e & 尾数 f & 值 & 十六进制例 \\
\midrule
+0 & 全 0 & 全 0 & $+0$ & \code{0x00000000} \\
-0 & 全 0 & 全 0 & $-0$（s=1） & \code{0x80000000} \\
非规格化数 & 全 0 & 非 0 & $(-1)^{s} \times 0.f \times 2^{-126}$ & \code{0x00000001} 为 $2^{-149}$ \\
$\pm$Inf & 全 1 & 全 0 & $(-1)^{s} \times \infty$ & \code{0x7F800000} 为 $+\infty$ \\
NaN & 全 1 & 非 0 & 不是一个数 & \code{0x7FC00000} \\
\bottomrule
\end{longtable}

非规格化数处理下溢一侧的边界。规格化数依赖隐含的前导 1，能表示的最小正数是 \code{e = 1}、\code{f = 0} 的 $1.0 \times 2^{-126}$（\code{0x00800000}）；比这更小的数装不下隐含位。于是标准约定 \code{e = 0} 时隐含位变成 0，指数固定为 -126，值为 $0.f \times 2^{-126}$。f 的 23 个 bit 里最小非零值是末位的 1，对应 $2^{-23} \times 2^{-126} = 2^{-149}$，这就是单精度最小正数。关键在于过渡方式：从 $2^{-126}$ 继续向下，表示能力不是悬崖式掉到 0，而是有效位一位一位减少、精度逐渐丢失——这称为渐进下溢（gradual underflow）。典型产生场景是两个接近的极小规格化数相减，差落在规格化范围之下。

无穷大覆盖另一端的边界，产生路径有两条。一条是真除零：IEEE 语义下 $1.0 / 0.0$ 得 $+\infty$、$-1.0 / 0.0$ 得 $-\infty$，同时置除零标志，默认不停机。另一条是溢出：结果的指数超过单精度上限（约 $3.4 \times 10^{38}$），例如 $10^{38} \times 10$，按当前舍入模式变成 Inf 或最大有限值。Inf 参与运算有定义好的结果：$1 / \infty = 0$，$\infty + x = \infty$，比较时大于一切有限值。流水线因此不必为超界结果插入特例停顿。

NaN 表示“没有定义”的运算结果：$0/0$、$\infty - \infty$、$0 \times \infty$、对负数开平方都产生 NaN。硬件必须实现它的两条语义。一是传播：NaN 参与几乎所有运算，结果仍是 NaN（尾数字段通常原样传下去），异常不会在长算式中途被静默吞掉。二是比较：NaN 与任何值比较都为不相等，包括它自己，所以 \code{x == x} 为假即可检出 NaN——这条性质被编译器和数学库直接利用。还要注意零的符号：+0 与 -0 比较相等，但 $1/(+0) = +\infty$、$1/(-0) = -\infty$，符号位在除法中仍然可观察；负数下溢向零舍入时也会产生 -0。

特殊值的意义在于让浮点成为封闭系统：任意 32-bit 输入、任意运算，都有定义好的 32-bit 输出，控制逻辑无需为异常组合停下流水线。代价是 FPU 必须为阶码全 0 和全 1 各准备一条旁路，译码、比较、类型转换和舍入都要先查这两类模式——这正是前面“特殊值需要旁路和异常规则”的具体内容。

浮点加法比整数加法复杂得多。两个整数相加，只需对齐同名 bit 并沿进位链传播；两个浮点数相加，通常要先比较指数，把较小数的尾数右移对齐，再做加减，随后规格化结果、舍入，并检查异常标志。若一个很大的数和一个很小的数相加，小数可能在对阶右移时被移出保留位宽之外，结果看起来像“小数被吃掉”。这不是软件显示问题，而是 FPU 内部位宽和舍入规则的结果。

\begin{lstlisting}
large = 1.0 * 2^30
small = 1.0
large + small may round back to large
when the significand has too few bits.
\end{lstlisting}

\topic{一次浮点加法要走完对阶、相加、规格化、舍入四步}
FPU 的加法器不是一条进位链，而是一条固定次序的小流水线。用 $1.5 + 0.078125$ 把四步完整走一遍。先把两个十进制数写成规格化二进制，再写成 binary32：

\begin{lstlisting}
1.5      = 1.1_2  * 2^0
           s=0  e=0111 1111 (127)  f=100 0000 ... 0000
           -> 0x3FC00000
0.078125 = 1.01_2 * 2^-4
           s=0  e=0111 1011 (123)  f=010 0000 ... 0000
           -> 0x3DA00000
\end{lstlisting}

第一步，对阶。两数指数不同，尾数不能直接相加：小数点位置不同的两个值必须先对齐，这与前面定点加法的前提完全一样。指数差是 $127 - 123 = 4$，规则是小阶向大阶对齐：把 0.078125 的尾数右移 4 位，它的指数随之补到 127。为什么不让大阶向小阶对齐？那需要把大数尾数左移，最高的隐含位会移出保留位宽，丢的是最高有效位；右移小数丢的是最低有效位，损失最小。对阶后的 bit：

\begin{lstlisting}
1.5:       1.1000 0000 0000 0000 0000 0000   * 2^0
0.078125:  0.0001 0100 0000 0000 0000 0000   * 2^0
           (significand shifted right by 4)
\end{lstlisting}

实际硬件在尾数右侧还会多留 guard、round、sticky 三种附加位，接住被移出的 bit，供最后一步舍入查询。

第二步，尾数相加。对齐后的两个 24-bit 尾数（含隐含位）走普通整数进位链：

\begin{lstlisting}
  1.1000 0000 0000 0000 0000 0000
+ 0.0001 0100 0000 0000 0000 0000
= 1.1001 0100 0000 0000 0000 0000
\end{lstlisting}

第三步，规格化。和已经是 $1.x$ 形态，前导 1 在位，无需移动。（若和大于等于 2，例如 $1.5 + 1.5 = 11.0_2$，要右移 1 位并把指数加 1；若和小于 1，例如两个相近数相减，要左移直到前导 1 重新出现，指数相应减小。）

第四步，舍入。本例 23 位尾数恰好装下 \code{1001 0100 ...}，被舍弃部分为 0，不发生舍入，结果直接拼回三段：

\begin{lstlisting}
s=0  e=0111 1111  f=100 1010 0000 ... 0000
-> 0x3FCA0000 = 1.578125 = 1.5 + 0.078125
\end{lstlisting}

把四步连成一张数据流图，每一步的输入和输出都是明确的字段，次序固定。

\begin{pdfdiagram}
  \node[hw label] (ina) at (-0.1,1.35) {A（s, e, f）};
  \node[hw label] (inb) at (1.1,1.35) {B（s, e, f）};
  \node[hw compute, minimum width=2.45cm, align=center] (s1) at (0.5,0) {对阶\\指数比较，小阶尾数右移};
  \node[hw compute, minimum width=2.45cm, align=center] (s2) at (3.35,0) {尾数相加\\24-bit 进位链};
  \node[hw compute, minimum width=2.45cm, align=center] (s3) at (6.2,0) {规格化\\左规或右规，调整指数};
  \node[hw compute, minimum width=2.45cm, align=center] (s4) at (9.05,0) {舍入并拼回\\nearest even};
  \draw[hw bus] (ina.south) -- (ina.south |- s1.north);
  \draw[hw bus] (inb.south) -- (inb.south |- s1.north);
  \draw[hw bus] (s1) -- (s2);
  \draw[hw bus] (s2) -- (s3);
  \draw[hw bus] (s3) -- (s4);
  \node[hw label] at (0.5,-1.2) {对齐尾数，统一指数};
  \node[hw label] at (3.35,-1.2) {和（含 GRS 附加位）};
  \node[hw label] at (6.2,-1.2) {$1.f$ 形态，新指数};
  \node[hw label] at (9.05,-1.2) {结果 \code{0x3FCA0000}};
\end{pdfdiagram}
\diagramnote{浮点加法四步的输入输出：对阶按指数差把小阶尾数右移，移出位进入 guard、round、sticky 附加位；尾数相加走普通整数进位链；规格化把结果移回 $1.f$ 形态并同步加减指数；舍入按 round to nearest even 处理附加位后拼回 s/e/f 三段。四步次序固定，真实 FPU 把它们分到不同流水级。}

舍入规则换一个必然触发它的例子看：$2^{24} + 1$。$2^{24}$ 的尾数是 \code{1.000 ... 000}，指数 24；$1$ 对阶时要右移 24 位，它唯一的那个 1 被移到 23 位尾数之外。此时和的真值是 $2^{24} + 1$，恰好落在两个可表示值 $2^{24}$ 和 $2^{24} + 2$ 的正中间——这个指数下尾数末位的权重（ulp）是 2。默认舍入模式是 round to nearest even：就近舍入，平局时取尾数末位为偶数的那个。$2^{24}$ 的尾数末位是 0，为偶，于是结果舍回 $2^{24}$：这次加法真实执行了，加上的 1 却被格式位宽吃掉。若加的是 3，$2^{24} + 3$ 恰好落在 $2^{24} + 2$ 和 $2^{24} + 4$ 的正中间，到两边距离都是 1，又是一个平局：按取偶规则，$2^{24} + 4$ 的尾数末位是 0，为偶，于是结果舍到 $2^{24} + 4$。

这个机制解释了为什么大数加小数容易“吃掉”小数：尾数只有 23 位，ulp 随指数放大，小数的全部有效位可能都在 ulp 之下，对阶右移后被舍入规则抹平。这也意味着浮点加法的每一步都可能按当时的指数重新舍入。

由此也能理解浮点加法通常不满足严格结合律。$(a+b)+c$ 和 $a+(b+c)$ 可能在不同中间步骤舍入，最终得到不同的最低几位。工程上不能把浮点数当成实数域的无限精确元素；必须知道格式位宽、舍入模式、异常处理和比较策略。

\begin{chipcase}
8087 x87 协处理器把浮点执行从 8086 主整数路径旁边分离出来，说明浮点硬件从一开始就有明显的复杂度边界。后来 x86 通过 SSE、AVX 等 SIMD（单指令多数据）浮点/向量扩展，把多个浮点操作组织成并行数据通路；现代高性能 CPU 还会用流水线、转发、乱序调度和宽向量单元降低吞吐延迟。相反，许多小型 MCU 为了面积、功耗和确定性，仍然选择没有 FPU 或只提供较窄 FPU，由软件库、整数或定点路径承担一部分数值任务。
\end{chipcase}

\topic{数值格式决定硬件代价和延迟}同样是“加法”，整数、定点、BCD 和浮点的硬件路径并不相同。整数加法的核心是位对齐后的进位链；定点加法在格式一致时复用整数加法器；BCD 加法需要半字节合法性检测和十进制校正；浮点加法则增加对阶、尾数运算、规格化、舍入和特殊值处理。格式越灵活，硬件要承担的解释工作越多。

四种格式的取舍也因此不同。无符号/补码整数靠加法器、移位器和比较标志工作，简洁、快速、面积低，代价是范围固定、不能直接表示小数；定点用整数 ALU 加缩放、舍入和饱和，实时性好、适合 DSP 和控制场景，代价是尺度要人工管理；BCD 用二进制加法加十进制校正，十进制输入输出直接，代价是位宽利用率低、校正逻辑额外增加；浮点靠对阶、尾数、规格化、舍入和异常处理工作，动态范围大、通用科学计算方便，代价是面积、功耗、验证和延迟全面上升。

现代芯片能把很多数值运算做得很快，不是因为物理定律允许“没有延迟”，而是因为工程上同时用了多种手段。第一，常见路径被专门做成短关键路径，例如整数加法器采用超前进位或前缀结构，避免所有进位逐位等待。第二，复杂运算被流水线化，单次结果仍有若干周期延迟，但每个周期都能接收新操作，提高吞吐。第三，SIMD/向量单元把多个相同格式的数据并排处理，摊薄取指、译码和控制成本。第四，缓存、寄存器重命名、旁路转发和调度逻辑尽量让数据在核心附近流动，减少等待内存和长线传输的时间。

因此，“低延迟”必须分清两种说法：\hwkey{单次操作延迟}指一个结果从输入到可用需要多久，\hwkey{吞吐率}指稳定状态下每个周期能完成多少个操作。现代浮点乘加单元可能有多个周期的单次延迟，但由于流水线很深，可以在每个周期发射新的乘加；整数 ALU 可能单次延迟更短，但若数据依赖形成长链，仍会受到前一条结果可用时间限制。读芯片资料时，应同时看 latency、throughput、数据格式、向量宽度、异常模式和是否支持饱和/舍入。

移位器也是算术部件。左移一位通常相当于乘以 2，但前提是移出去的 bit 没有携带需要保留的信息。右移更容易出差别。逻辑右移在高位补 0；算术右移保留符号位，用于补码有符号数。

\begin{lstlisting}
logical right shift:    1000 >> 1 = 0100
arithmetic right shift: 1000 >> 1 = 1100
\end{lstlisting}

若 \code{1000} 按 4-bit 补码解释为 -8，算术右移得到 \code{1100}，也就是 -4；逻辑右移得到 \code{0100}，也就是 4。两个结果差异巨大。原因不是移位器“出错”，而是控制信号选择了不同硬件规则。循环移位又是第三种规则：移出的 bit 从另一端回填，并且常常与进位标志一起参与多字长位操作。

\begin{chipcase}
8086 的 FLAGS 寄存器把 Carry、Zero、Sign、Overflow 等标志暴露给分支、条件转移和多字长运算。80386 把通用寄存器和地址计算扩展到 32-bit 后，仍保留同类标志语义。这个延续说明：位宽可以扩大，执行单元可以复杂化，但 CPU 仍需要一组可复查的\hwevidence{结果标志}，把 ALU 的组合输出交给控制流和后续指令使用。
\end{chipcase}

若想把本节内容搭到面包板上，可以先做一个 4-bit 数值解释实验：用 4 个拨码开关给出 \code{b3..b0}，用 4 个 LED 显示原始 bit，再用纸面表格分别写出无符号值和补码值。此时不需要任何复杂芯片，关键是训练一个习惯：LED 只显示 bit，数值来自解释规则。后面接入 74HC283 加法器、74HC86 反相控制和 74HC157 选择器时，仍然要保留这张解释表，否则很容易把结果 bit 与数学整数混为一谈。

\section{二、从加法器到可复用 ALU}

ALU 是算术逻辑单元。加法、减法、按位逻辑、移位和比较反复出现，硬件需要把这些能力集中到一个受控制的组合逻辑块中。ALU 的本质是：多个候选计算电路并排存在，控制信号选择本周期采用的结果，并同时产生可供控制器判断的标志位。

最小 ALU 可以写成接口：

\begin{lstlisting}
ALU(a, b, op) -> result, flags
\end{lstlisting}

\code{a} 和 \code{b} 是两个操作数，\code{op} 是选择信号，\code{result} 是结果，\code{flags} 是结果附带的标志信息。这里的接口不是软件函数，而是一组并行电线：两个操作数总线进入 ALU，控制线进入选择网络，结果总线和标志线从 ALU 输出。

\begin{longtable}{L{0.18\linewidth}L{0.32\linewidth}L{0.38\linewidth}}
\toprule
op & 运算 & 硬件来源 \\
\midrule
ADD & \code{a + b} & 加法器 \\
SUB & \code{a - b} & 加法器、B 反相控制、最低位进位 \\
AND & \code{a \& b} & 按位 AND 门阵列 \\
OR & \code{a | b} & 按位 OR 门阵列 \\
XOR & \code{a xor b} & 按位 XOR 门阵列 \\
SHIFT & 逻辑/算术/循环移位 & 移位器和补位选择逻辑 \\
CMP & 比较 & 通常复用减法路径和标志位 \\
\bottomrule
\end{longtable}

加法器从半加器开始。半加器接受两个 1-bit 输入，输出和位 \code{sum} 与进位 \code{carry}：

\begin{lstlisting}
sum   = a xor b
carry = a and b
\end{lstlisting}

全加器（完整加法器）还要接受来自低位的进位 \code{cin}：

\begin{lstlisting}
sum  = a xor b xor cin
cout = (a and b) or (cin and (a xor b))
\end{lstlisting}

把多个全加器串起来，就得到串行进位加法器。最低位先计算，再把进位送到下一位。这个结构规则、容易画出、容易搭电路，也最能说明关键路径：最高位结果必须等待低位进位逐级传播。位宽越宽，最坏延迟越大。

半加器用 XOR 和 AND 解释一位相加的和位与进位；全加器在半加器上加入进位合成，构成多位串行进位链；四个全加器级联就是一片 4-bit 加法器，对应 74HC283/74LS283 一类芯片；ALU 加法路径则再加上输入选择和标志生成，支持 ADD、SUB、地址计算和比较。

ALU 里最值得细看的是 SUB。若单独做减法器，会增加面积和连线。补码让它复用 ADD 的加法器：

\begin{lstlisting}
if op = ADD:
  b_to_adder = b
  carry_in   = 0

if op = SUB:
  b_to_adder = NOT b
  carry_in   = 1

result = a + b_to_adder + carry_in
\end{lstlisting}

实际电路里，B 的每一位前面可以接 XOR 门：当 \code{sub=0}，\code{b xor sub = b}；当 \code{sub=1}，\code{b xor sub = NOT b}。最低位进位输入也由 \code{sub} 控制。这样 ADD 和 SUB 共用同一条进位链。

\begin{longtable}{L{0.18\linewidth}L{0.24\linewidth}L{0.24\linewidth}L{0.24\linewidth}}
\toprule
\code{sub} & B 输入处理 & \code{Cin0} & 实际运算 \\
\midrule
0 & B 原样通过 & 0 & \code{A + B} \\
1 & B 每位取反 & 1 & \code{A + (~B) + 1 = A - B} \\
\bottomrule
\end{longtable}

这个小结构可以直接搭成可观察实验。实验不需要先做完整 CPU，只要把几个连接点分开搭建、逐一观察，就已经包含了 ALU 复用的核心：A/B 输入用拨码开关给出两个 4-bit 操作数，原始 bit 要和纸面数值解释一致；B 反相用 74HC86 的四个 XOR 门处理 B 的四位，\code{sub=0} 时原样通过、\code{sub=1} 时逐位取反；进位输入把同一个 \code{sub} 接到最低位 \code{Cin}，加法时 \code{Cin=0}、减法时 \code{Cin=1}；结果输出让 74HC283 的输出接 LED 或逻辑探针，\code{sub} 切换时加法和减法结果都要可复现。

\begin{chipcase}
74HC283/74LS283 是常见 4-bit 二进制全加器，适合观察多位进位链；74HC86 提供四个二输入 XOR 门，适合做 B 输入的可控反相；74HC157 提供四路二选一选择器，可用于在逻辑结果和加法结果之间选择。若再接上 74HC173 或 74HC574 保存操作数和结果，就能组成一个可单步观察的 4-bit ALU 实验台。实验重点不是追求速度，而是把 \code{sub}、\code{Cin}、结果 LED 和标志 LED 之间的关系看清楚。
\end{chipcase}

ALU 内部可以并行产生多个候选结果，再用多路选择器按 \code{op} 选出最终 \code{result}。

\begin{lstlisting}
add_result   = adder(a, b_to_adder, carry_in)
and_result   = a and b
or_result    = a or b
xor_result   = a xor b
shift_result = shifter(a, shift_amount, shift_kind)

result = mux(op, add_result, and_result,
             or_result, xor_result, shift_result)
\end{lstlisting}

这不表示所有运算按软件顺序“算一遍”。组合逻辑是并行传播的：加法器、按位逻辑、移位器都可能同时产生自己的候选输出。最终哪一路被后级看见，由 mux 决定。代价也随之出现。候选电路越多，ALU 面积越大；mux 输入越多，选择延迟越大；若某个候选结果来自长进位链，它可能成为 ALU 的关键路径。

ALU 常见标志位包括：

\begin{longtable}{L{0.18\linewidth}L{0.36\linewidth}L{0.36\linewidth}}
\toprule
标志位 & 来源 & 用途 \\
\midrule
Zero & 结果所有 bit 做 OR 后再取反 & 相等判断、循环结束、条件选择 \\
Carry & 最高位外进位或借位约定 & 无符号溢出、多字长加减、无符号比较 \\
Overflow & 补码有符号运算的符号异常 & 有符号范围检查、有符号比较 \\
Negative/Sign & 结果最高位 & 补码符号判断、条件分支输入 \\
Parity/辅助标志 & 特定低位统计或半字节进位 & 某些 ISA、BCD 或兼容性需求 \\
\bottomrule
\end{longtable}

Zero 可以由所有结果 bit 做 OR 后再取反得到。Negative 通常直接来自结果最高位。Carry 来自加法器最高位外的进位。Overflow 则要看有符号加减法中符号是否出现不可能的变化。对于加法，若两个输入符号相同而结果符号不同，则发生补码溢出；对于减法，若被减数和减数符号不同而结果符号与被减数不同，也发生补码溢出。

四个标志的生成电路各不相同，但都与结果同时产生：它们并联在结果总线和进位链上，互不等待。

\begin{pdfdiagram}
  \node[hw state, minimum width=2.6cm, align=center] (res) at (0,0.8) {结果总线 R[7:0]};
  \node[hw compute, minimum width=2.6cm, align=center] (chain) at (0,-1.2) {加法器进位链};
  \node[hw compute, minimum width=2.2cm, align=center] (fz) at (5.3,1.6) {NOR\\（全部 bit）};
  \node[hw compute, minimum width=2.2cm, align=center] (fn) at (5.3,0.4) {取最高位 R7};
  \node[hw compute, minimum width=2.2cm, align=center] (fc) at (5.3,-1.2) {末级进位};
  \node[hw compute, minimum width=2.2cm, align=center] (fv) at (5.3,-2.6) {Cout XOR C7};
  \draw[draw=BookMuted, line width=0.62pt] (res.east) -- (2.4,0.8);
  \draw[draw=BookMuted, line width=0.62pt] (2.4,1.6) -- (2.4,0.4);
  \draw[hw bus] (2.4,1.6) -- (fz.west);
  \node[hw label] at (3.3,1.78) {8 bit};
  \draw[hw bus] (2.4,0.4) -- (fn.west);
  \node[hw label] at (3.3,0.58) {R7};
  \draw[draw=BookMuted, line width=0.62pt] (chain.east) -- (2.4,-1.2);
  \draw[draw=BookMuted, line width=0.62pt] (2.4,-1.2) -- (2.4,-2.6);
  \draw[hw bus] (2.4,-1.2) -- (fc.west);
  \node[hw label] at (3.3,-1.02) {Cout};
  \draw[hw bus] (2.4,-2.6) -- (fv.west);
  \node[hw label] at (3.3,-2.42) {Cout、C7};
  \draw[hw bus] (fz.east) -- ++(1.0,0) node[right, hw label] {Z};
  \draw[hw bus] (fn.east) -- ++(1.0,0) node[right, hw label] {N};
  \draw[hw bus] (fc.east) -- ++(1.0,0) node[right, hw label] {C};
  \draw[hw bus] (fv.east) -- ++(1.0,0) node[right, hw label] {V};
\end{pdfdiagram}
\diagramnote{四个标志并联生成，与结果同时有效：Z 是结果全部 bit 的 NOR；N 直接取结果最高位 R7；C 是进位链末级输出 Cout；V 用末两级进位 Cout XOR C7——最高位与次高位进位不一致即补码溢出。四根标志线随后送进标志寄存器，供条件分支查询。}

\begin{longtable}{L{0.16\linewidth}L{0.28\linewidth}L{0.28\linewidth}L{0.18\linewidth}}
\toprule
动作 & Carry/Borrow 解释 & Overflow 解释 & 常见用途 \\
\midrule
ADD & 最高位外进位，表示无符号超界 & 同号相加得到异号 & 无符号加法、多字长加法、有符号溢出检查 \\
SUB & 由 ISA 约定表示无借位或借位 & 异号相减且结果符号异常 & 无符号大小比较、有符号大小比较 \\
CMP & 通常执行减法但不写回结果 & 与 SUB 相同，只保留比较标志 & 条件分支、条件选择、循环判断 \\
INC/DEC & 有些 ISA 不更新 Carry，或按专门规则更新 & 仍按补码边界判断 & 地址递增、循环计数，必须查清 ISA 约定 \\
\bottomrule
\end{longtable}

这张表强调一个容易被忽略的事实：标志位不是全局数学真理，而是\hwevidence{指令语义与硬件标志的约定}。同样执行 \code{A-B}，相等判断只关心 Zero，无符号小于关心借位或 Carry 约定，有符号小于要结合 Sign 和 Overflow；判断 A 是否为 0，直接检查 A 或执行传递操作，看 Zero 标志即可。早期 x86、6502、Z80 等处理器都把这些标志放进状态寄存器，但具体命名、是否更新、条件码如何解释并不完全相同。读芯片手册或写控制器时，必须先弄清“本条动作会更新哪些标志”。

多字长运算展示了 Carry 的实用价值。若只有 8-bit ALU，却要做 16-bit 加法，可以先加低 8 位，保存低字节结果和进位；再加高 8 位，并把低字节产生的 Carry 作为高字节加法的 \code{Cin}。早期 8-bit 处理器经常依靠这种方式完成更宽整数运算。软件看到的是 16-bit 加法，硬件实际执行的是多个周期内的低字节、高字节和进位链传递。

\begin{lstlisting}
low_sum, carry0  = add8(A_low,  B_low,  0)
high_sum, carry1 = add8(A_high, B_high, carry0)
result = high_sum:low_sum
\end{lstlisting}

把这个观察放到经典芯片中，可以看到同一问题的不同规模。6502、Z80、8086、80386 都必须解决 ALU 如何支持加法、减法、逻辑运算、移位和标志位。早期 8-bit 处理器更强调用较小数据通路和多周期控制完成动作；8086 把 16-bit 数据通路、段偏移地址形成和标志控制组织在更复杂的执行部件里，ALU 不只做算术，也服务地址计算和条件判断；80386 扩展到 32-bit 后，位宽、标志、地址计算和控制路径都随之加重，进位、布线、译码和时序压力同时上升；现代核心则把 ALU 复制到多个执行单元，用旁路、重命名和乱序调度喂养它们，但每条执行路径仍要满足时序约束。不同芯片的实现细节不同，但抽象不变：候选计算、选择器、进位路径、标志生成和时钟边界共同形成 ALU 行为。

\topic{超前进位按层级展开，进位选择用复制换等待}
串行进位的瓶颈前面已经看过：16-bit 加法器的 \code{C16} 要等 16 段进位接力。超前进位（carry-lookahead）把每一位的进位方程 $C_{i+1} = G_i + P_i C_i$ 逐位代入，直到右边只剩 $G$、$P$ 和 \code{C0}。4-bit 完整展开：

\begin{lstlisting}
G_i = a_i AND b_i        (generate)
P_i = a_i XOR b_i        (propagate)

C1 = G0 + P0*C0
C2 = G1 + P1*G0 + P1*P0*C0
C3 = G2 + P2*G1 + P2*P1*G0 + P2*P1*P0*C0
C4 = G3 + P3*G2 + P3*P2*G1 + P3*P2*P1*G0 + P3*P2*P1*P0*C0
\end{lstlisting}

右边只依赖输入位和 \code{C0}，四个进位因此可以同时开算：一串代入被换成一片与-或两级逻辑。代价也写在式子里——项数和扇入随位宽增长，\code{C4} 的最后一项要 5 输入 AND，进位 OR 要收 4 项；位宽到 8 以上，单级展开的门扇入就不现实。所以更宽的加法器把同一条递推分层使用。

16-bit 的做法是 4 组 4-bit CLA 再加一层组间逻辑。每组内部按上面的展开并行产生组内进位；每组再输出两个组级信号：组产生 \code{GG} 表示“本组不管低位进位是什么，自己一定向外产生进位”，组传播 \code{GP} 表示“低位来什么进位，本组原样传出去”：

\begin{lstlisting}
GG = G3 + P3*G2 + P3*P2*G1 + P3*P2*P1*G0
GP = P3*P2*P1*P0
\end{lstlisting}

第二级用完全相同的递推处理四个组的 \code{GG}/\code{GP}，一次算出全部组间进位：

\begin{lstlisting}
C4  = GG0 + GP0*C0
C8  = GG1 + GP1*GG0 + GP1*GP0*C0
C12 = GG2 + GP2*GG1 + GP2*GP1*GG0 + GP2*GP1*GP0*C0
C16 = GG3 + GP3*GG2 + GP3*GP2*GG1 + GP3*GP2*GP1*GG0
      + GP3*GP2*GP1*GP0*C0
\end{lstlisting}

延迟由此重排。串行 16-bit 是 16 段进位接力；两级超前的进位路径只剩约 5 段：位级 G/P 一级门，组级 \code{GG}/\code{GP} 生成与-或两级，组间进位与-或两级，组内进位与-或两级，末位求和 XOR 一级。按级门延迟估算，串行约 $1 + 2 \times 16 = 33$ 级，两级超前约 $1+2+2+2+1=8$ 级——接力段数从 16 降到约 5。层级化的含义就是：$C = G + P\cdot C$ 这条递推在组内用一次、在组间再用一次；经典板级方案里，74182 超前进位发生器承担的正是组间这一层。

把 16-bit 的两级结构画成方块图，每个组对外的接口只剩 GG、GP 和组间进位三根线。

\begin{pdfdiagram}
  \node[hw compute, minimum width=2.3cm, minimum height=1.1cm, align=center] (g0) at (0,0) {4-bit CLA 组 0\\\code{Cin}=C0};
  \node[hw compute, minimum width=2.3cm, minimum height=1.1cm, align=center] (g1) at (2.55,0) {4-bit CLA 组 1};
  \node[hw compute, minimum width=2.3cm, minimum height=1.1cm, align=center] (g2) at (5.1,0) {4-bit CLA 组 2};
  \node[hw compute, minimum width=2.3cm, minimum height=1.1cm, align=center] (g3) at (7.65,0) {4-bit CLA 组 3};
  \node[hw control, minimum width=8.8cm, minimum height=1.1cm, align=center] (cla2) at (3.9,2.1) {组间超前进位逻辑（第二级）：C4、C8、C12、C16 并行产生};
  \node[hw label] (c0) at (-1.5,2.1) {C0};
  \draw[hw bus] (c0.east) -- (cla2.west);
  \draw[hw bus] (0,0.55) -- (0,1.55);
  \node[hw label, anchor=east, align=center] at (-0.15,1.05) {GG0\\GP0};
  \draw[hw bus] (2.05,0.55) -- (2.05,1.55);
  \node[hw label, anchor=east, align=center] at (1.9,1.05) {GG1\\GP1};
  \draw[hw bus] (4.6,0.55) -- (4.6,1.55);
  \node[hw label, anchor=east, align=center] at (4.45,1.05) {GG2\\GP2};
  \draw[hw bus] (7.15,0.55) -- (7.15,1.55);
  \node[hw label, anchor=east, align=center] at (7.0,1.05) {GG3\\GP3};
  \draw[hw bus] (3.1,1.55) -- (3.1,0.55);
  \node[hw label, anchor=west] at (3.25,1.05) {C4};
  \draw[hw bus] (5.65,1.55) -- (5.65,0.55);
  \node[hw label, anchor=west] at (5.8,1.05) {C8};
  \draw[hw bus] (8.2,1.55) -- (8.2,0.55);
  \node[hw label, anchor=west] at (8.35,1.05) {C12};
  \draw[hw bus] (cla2.east) -- ++(0.7,0) node[above, hw label] {C16};
  \draw[hw bus] (g0.south) -- ++(0,-0.65) node[below, hw label] {S[3:0]};
  \draw[hw bus] (g1.south) -- ++(0,-0.65) node[below, hw label] {S[7:4]};
  \draw[hw bus] (g2.south) -- ++(0,-0.65) node[below, hw label] {S[11:8]};
  \draw[hw bus] (g3.south) -- ++(0,-0.65) node[below, hw label] {S[15:12]};
\end{pdfdiagram}
\diagramnote{16-bit 两级超前进位：四个 4-bit CLA 组各自生成组产生 GG 与组传播 GP 上报第二级；第二级用同一条 $C=G+P\cdot C$ 递推并行算出 C4、C8、C12、C16 发回各组（经典 74182 承担的正是这一层）。进位接力从串行 16 段降到约 5 段，约 8 级门。}

进位选择加法器（carry-select）走另一条路。把 16-bit 分成 4 段，每段并排摆两套 4-bit 加法器：一套假设段进位输入为 0，另一套假设为 1，两份和同时算完。真实的段间进位到达时，不再穿过段内进位链，只驱动一层 mux，从两份候选里选一份。等待时间变成“第一段完整加法”加“其后每段一级 mux”，约等于 4-bit 加法加三级 mux，远小于 16 段接力。

代价同样直接：段内加法器接近两倍，再加一层与位宽成正比的 mux 阵列，面积显著上升。为什么仍值得？因为串行结构里“等进位”的那段时间被利用起来，把两种可能的分支都提前算完——选择代替了等待。\hwkey{提前算出所有候选，再用真实信号选择}是硬件里反复出现的换时手段：超前进位用逻辑深度换时间，进位选择用电路复制换时间；位宽继续加大时两者常常组合使用（组内超前、组间选择），32-bit 和 64-bit ALU 的主加法器很少是单一结构。

\begin{longtable}{L{0.15\linewidth}L{0.28\linewidth}L{0.27\linewidth}L{0.2\linewidth}}
\toprule
16-bit 结构 & 进位等待 & 额外代价 & 适用场景 \\
\midrule
串行进位 & 16 段接力，约 33 级门 & 几乎无额外面积 & 教学实验、窄位宽 \\
两级超前 & 约 5 段，约 8 级门 & 大扇入与-或逻辑、组间 PG 网络 & ALU 主加法路径 \\
进位选择 & 一段加法加每段一级 mux & 段内加法器近两倍加 mux 阵列 & 高主频宽位加法 \\
\bottomrule
\end{longtable}

加法器的位宽不是简单复制全加器，进位组织方式本身就是设计变量。

\section{三、进位、移位、比较与乘法的硬件取舍}

加法器是 ALU 的核心路径之一，因为减法、地址计算、数组下标、栈指针更新和分支目标计算都会反复使用它。最直观的结构是串行进位加法器：第 0 位产生的进位送到第 1 位，第 1 位再送到第 2 位，直到最高位。它面积小、结构规则，但最坏延迟随位宽增加而增长。4-bit 时容易接受，32-bit 或 64-bit 时就可能成为关键路径。

可以把每一位的进位逻辑归结为两类信号：

\begin{lstlisting}
generate_i  = a_i and b_i
propagate_i = a_i xor b_i
carry_out_i = generate_i or (propagate_i and carry_in_i)
\end{lstlisting}

若本位的 A 和 B 都为 1，本位会\hwkey{产生}进位；若 A 和 B 有且仅有一个为 1，本位会\hwtiming{传递}低位进位。快速加法器的思想，就是用额外组合逻辑提前计算这些产生和传递关系，减少“等低位一位一位传上来”的时间。这说明速度不是免费获得的：更短的延迟通常换来更多门、更复杂连线和更严格布局。

把“提前计算”真正展开一次。以 4-bit 为例，把 $C_{i+1} = G_i \lor (P_i C_i)$ 逐位代入，每个进位都改写成只含 G、P 和 \code{C0} 的形式——这正是第二节已经展开过的四行式子。关键观察不变：右边只依赖 A、B 和 \code{C0}，\code{C4} 不再等 \code{C3}，四个进位可以同时开算。这就是“超前进位”的全部含义——不是让进位消失，而是把进位链的串行代入改成并行的与或逻辑。

\begin{pdfdiagram}
  \node[hw compute, minimum width=1.6cm, minimum height=0.85cm] (gp0) at (0,0) {G0, P0};
  \node[hw compute, minimum width=1.6cm, minimum height=0.85cm] (gp1) at (2.4,0) {G1, P1};
  \node[hw compute, minimum width=1.6cm, minimum height=0.85cm] (gp2) at (4.8,0) {G2, P2};
  \node[hw compute, minimum width=1.6cm, minimum height=0.85cm] (gp3) at (7.2,0) {G3, P3};
  \node[hw label] at (0,0.95) {A0, B0};
  \node[hw label] at (2.4,0.95) {A1, B1};
  \node[hw label] at (4.8,0.95) {A2, B2};
  \node[hw label] at (7.2,0.95) {A3, B3};
  \draw[hw bus] (0,0.65) -- (gp0.north);
  \draw[hw bus] (2.4,0.65) -- (gp1.north);
  \draw[hw bus] (4.8,0.65) -- (gp2.north);
  \draw[hw bus] (7.2,0.65) -- (gp3.north);
  \node[hw control, minimum width=8.2cm, minimum height=1.1cm] (cla) at (3.6,-1.75) {进位逻辑（与-或两级）：C1、C2、C3、C4 并行产生};
  \draw[hw bus] (gp0.south) -- (0,-1.2);
  \draw[hw bus] (gp1.south) -- (2.4,-1.2);
  \draw[hw bus] (gp2.south) -- (4.8,-1.2);
  \draw[hw bus] (gp3.south) -- (7.2,-1.2);
  \node[hw label] at (-1.6,-1.75) {\code{C0}};
  \draw[hw bus] (-1.2,-1.75) -- (cla.west);
  \draw[hw bus] (cla.east) -- ++(0.5,0) node[right, hw label]{C1};
  \draw[hw bus] (cla.east |- cla.north east) ++(0,0) -- ++(0.5,0) node[right, hw label]{C2};
  \draw[hw bus] (cla.east |- cla.south east) ++(0,0) -- ++(0.5,0) node[right, hw label]{C3};
  \draw[hw bus] (cla.south) -- ++(0,-0.45) node[below, hw label]{C4};
\end{pdfdiagram}
\diagramnote{4-bit 超前进位结构：四个 G/P 单元同时算出产生/传播信号，进位逻辑用与-或两级把 C1–C4 一次性算出来，不再串行等待。}

延迟对比按“级门延迟”估算（数量级示意，不代表具体工艺）：4-bit 串行进位，G/P 1 级加每位约 2 级进位逻辑再乘 4 级接力，最坏约 9 级；4-bit 超前进位，G/P 1 级、进位逻辑与-或 2 级、求和 XOR 1 级，最坏约 4 级。位宽越大差距越夸张：64-bit 串行要约 $1+2\times64=129$ 级，而树形前缀结构能把进位逻辑压到 $\log_2 64=6$ 层左右。这就是 64-bit ALU 不能只按“64 个全加器排队”估算性能的原因——真实高性能加法器都在用前缀或分组结构买这几级延迟。

几种结构的取舍由此而来：串行进位面积小、布线规则、容易扩展，代价是进位逐位传播、最坏延迟随位宽增长；超前进位预先计算产生/传递关系、缩短进位等待，代价是组合逻辑和布线更复杂；分组进位在面积和速度之间折中，代价是要设计分组边界和跨组进位；前缀加法器适合高性能宽位加法，代价是布线密集、功耗和布局压力较大。教学机可以先用串行进位把 trace 跑通，追求主频时再换超前或前缀结构。

现代芯片为什么能降低延迟，不能简单归结为“晶体管更快”。在真实处理器里，低延迟来自一组共同措施：把常用路径做短，把宽路径拆分成前缀或分组结构，把长组合逻辑放进流水线，把结果通过旁路网络尽快送给后续执行单元，把缓存和寄存器堆放在物理上更近的位置，并让版图工具控制连线长度、扇出和时钟偏斜。晶体管速度、门级结构、物理布局、流水线边界和调度策略共同决定最终周期时间。每项措施都有代价：更快的标准单元可能增加面积、功耗和漏电；分组和前缀进位让布线更密、布局更难；流水线切分把延迟折成周期数，控制和旁路更复杂；结果旁路要增加选择网络和冲突判断；物理邻近要求模块复制和更强的布局约束；专用单元则要求常见运算有足够高的利用率才值得定制。

移位器也有类似取舍。若每次只移动一位，硬件简单，但一次移动 13 位就要多周期完成。桶形移位器允许在一个组合路径中按控制位选择移 1、2、4、8 等距离，多个阶段叠加得到任意移位量。以 8-bit 为例，移位量的低三位可以控制三层 mux：第一层决定是否移 1 位，第二层决定是否再移 2 位，第三层决定是否再移 4 位。这样能在一个周期内完成多位移位，但 mux 层数和连线负载会进入关键路径。

\begin{pdfdiagram}
  \node[hw state, minimum width=1.6cm] (din) at (0,0) {D[7:0]};
  \node[hw compute, minimum width=2.5cm] (m1) at (3.1,0) {8× 2:1 mux\\移 1 位？};
  \node[hw compute, minimum width=2.5cm] (m2) at (6.7,0) {8× 2:1 mux\\移 2 位？};
  \node[hw compute, minimum width=2.5cm] (m3) at (10.3,0) {8× 2:1 mux\\移 4 位？};
  \node[hw state, minimum width=1.6cm] (dout) at (13.3,0) {Y[7:0]};
  \draw[hw bus] (din) -- (m1);
  \draw[hw bus] (m1) -- (m2);
  \draw[hw bus] (m2) -- (m3);
  \draw[hw bus] (m3) -- (dout);
  \node[hw control, minimum width=1.1cm] (s0) at (3.1,-1.6) {s0};
  \node[hw control, minimum width=1.1cm] (s1) at (6.7,-1.6) {s1};
  \node[hw control, minimum width=1.1cm] (s2) at (10.3,-1.6) {s2};
  \draw[hw bus] (s0) -- (m1.south);
  \draw[hw bus] (s1) -- (m2.south);
  \draw[hw bus] (s2) -- (m3.south);
  \node[hw label, text width=11.5cm] at (6.4,-2.7) {例：移位量为 5（\code{101}）时，s2=1 先移 4 位，s1=0 不动，s0=1 再移 1 位，合计 5 位，一次组合传播完成。};
\end{pdfdiagram}
\diagramnote{8-bit 桶形移位器：三层 mux 分别由移位量的三个 bit 控制，每一层选择“直送”或“移 $2^k$ 位”。任意 0–7 位移位量都能在一个组合路径内完成，代价是三层 mux 的延迟和大量交叉连线。}

几种移位的硬件规则也要分清：逻辑移位在空出的一端补 0，控制线选错会改变数值解释；算术右移在高位补符号位，必须区分有符号和无符号解释；桶形移位用多层 mux 按 1/2/4/8 位选择，mux 延迟和长连线容易成为关键路径；循环移位把移出位从另一端回填，标志位和数据位的关系必须写清楚。

比较器通常可以复用减法路径。判断相等时，执行 \code{A-B} 后检查 Zero；判断无符号大小时，关注借位或 Carry 的解释；判断补码有符号大小时，不应仅看最高位，还要结合 Overflow。例如 \code{1000} 和 \code{0111} 在补码中分别是 -8 和 7，在无符号中分别是 8 和 7。同一组 bit 的大小关系随解释变化，因此比较器必须和控制器约定当前比较类型。

\begin{lstlisting}
eq_unsigned_or_signed = (A - B).Zero
lt_unsigned           = borrow_from_subtract
lt_signed             = result.Sign xor result.Overflow
\end{lstlisting}

这里的公式表达的是常见标志组合，具体处理器对 Carry 与 Borrow 的命名可能不同。重要的是不要把“最高位为 1”误当成所有比较场景里的“小于”。在补码有符号比较中，结果符号需要用 Overflow 修正；在无符号比较中，符号位没有负数含义。

乘法器展示了另一类工程折中。最朴素的二进制乘法可以生成部分积，再把部分积相加。若一次性展开成组合阵列，延迟和面积都较大，但吞吐可以高；若用移位加法多周期完成，面积小，但一次乘法需要多个周期。早期处理器常常用多周期或微码方式处理复杂运算，现代处理器则可能提供专用乘法单元、流水化乘法器或多个执行单元。抽象上看，乘法仍然是“部分积生成、对齐、求和、保存结果”，只是硬件把这些步骤压缩、展开或流水线化。

几种实现的取舍：移位加法每周期检查乘数一位、必要时累加被乘数移位值，面积小、周期数多，适合简单控制器；组合阵列同时生成多个部分积并通过加法网络求和，面积大、延迟大，但吞吐高；流水化乘法器在部分积压缩和求和之间插入寄存器，吞吐高但单次结果跨多个周期；Booth 编码和压缩树减少部分积数量或加快求和，代价是控制和布局复杂度增加。

\topic{阵列乘法器把乘法展开成部分积矩阵}

二进制乘法和十进制竖式乘法结构相同：用乘数的每一位去乘被乘数，把得到的行按位权错开，再全部相加。差别只在于二进制里“一位乘一个数”格外简单——乘数位是 1，部分积就是被乘数本身；乘数位是 0，部分积就是全 0。所以“生成一个部分积位”在硬件上就是一个 AND 门：$p_{ij}=a_i \land b_j$。

以 4×4 位无符号乘法 $A\times B$ 为例，$A=a_3a_2a_1a_0$，$B=b_3b_2b_1b_0$，先把 16 个部分积位排成矩阵，每一行是被乘数与乘数一位的逐位与，行与行之间按乘数位的权重错开一位：

\begin{lstlisting}
                a3   a2   a1   a0
             x  b3   b2   b1   b0
  ---------------------------------
                a3b0 a2b0 a1b0 a0b0    <- 第 0 行，权 2^0
           a3b1 a2b1 a1b1 a0b1         <- 第 1 行，权 2^1
      a3b2 a2b2 a1b2 a0b2              <- 第 2 行，权 2^2
 a3b3 a2b3 a1b3 a0b3                   <- 第 3 行，权 2^3
  ---------------------------------
 p7  p6  p5  p4  p3  p2  p1  p0        <- 乘积，共 8 位
\end{lstlisting}

代入具体数值 $0111\times0101$（$7\times5=35$），逐行移位累加：

\begin{lstlisting}
        0 1 1 1
     x  0 1 0 1
  -------------
        0 1 1 1      (b0=1，部分积 = 被乘数)
      0 0 0 0        (b1=0，部分积 = 0)
    0 1 1 1          (b2=1，左移两位)
  0 0 0 0            (b3=0)
  -------------
  0 0 1 0 0 0 1 1  = 35
\end{lstlisting}

阵列乘法器就是把这张竖式直接铺成硬件：$n^2$ 个 AND 门一次生成全部部分积位，再用 $n-1$ 行 n 位加法器逐行累加——第 0 行与第 1 行相加，结果再与第 2 行相加，依此类推。4×4 位乘法因此需要 16 个 AND 门和 3 行加法器，共约 12 个全加器。

这张竖式对应到硬件，就是下面的结构：AND 阵列一次生成 16 个部分积位，三行 4-bit 加法器逐行移位累加。

\begin{pdfdiagram}
  \node[hw label] (ma) at (1.8,1.85) {A[3:0]（被乘数）};
  \node[hw label] (mb) at (5.2,1.85) {B[3:0]（乘数）};
  \node[hw compute, minimum width=6.8cm, align=center] (and) at (3.5,0.75) {AND 阵列：16 个部分积位 $p_{ij}=a_i \land b_j$};
  \node[hw compute, minimum width=5.2cm, align=center] (ad1) at (3.5,-0.85) {加法器①：PP0 $+$ (PP1 $\ll 1$)};
  \node[hw compute, minimum width=5.2cm, align=center] (ad2) at (3.5,-2.25) {加法器②：和① $+$ (PP2 $\ll 2$)};
  \node[hw compute, minimum width=5.2cm, align=center] (ad3) at (3.5,-3.65) {加法器③：和② $+$ (PP3 $\ll 3$)};
  \draw[hw bus] (ma.south) -- (ma.south |- and.north);
  \draw[hw bus] (mb.south) -- (mb.south |- and.north);
  \draw[hw bus] (and.south) -- (ad1.north);
  \node[hw label, anchor=west] at (3.7,-0.05) {第 0、1 行部分积};
  \draw[hw bus] (ad1.south) -- (ad2.north);
  \node[hw label, anchor=west] at (3.7,-1.55) {部分和};
  \draw[hw bus] (ad2.south) -- (ad3.north);
  \node[hw label, anchor=west] at (3.7,-2.95) {部分和};
  \draw[hw bus] (ad3.south) -- ++(0,-0.7) node[below, hw label] {P[7:0] $= 0010\,0011_2 = 35$};
\end{pdfdiagram}
\diagramnote{4×4 阵列乘法器结构：16 个 AND 门一次生成全部部分积位，三行 4-bit 加法器逐行累加，每行按乘数位权重错开一位（$\ll 1$、$\ll 2$、$\ll 3$）。4×4 需要 16 个 AND 门和约 12 个全加器；$0111\times0101$ 经三行累加得 $0010\,0011_2=35$。}

延迟来自“逐行累加”这条链。每行加法本身有一次进位传播，行与行之间又是串行的，所以阵列乘法器的最坏延迟大致随位宽线性增长，且每一级都比同位宽加法器多出一行进位链。面积的增长更快：AND 门和全加器的数量都按 $n^2$ 走，连线规模也是 $n^2$ 量级。位宽翻一倍，加法器只长大一倍，乘法器却大约变为四倍。

\begin{longtable}{L{0.14\linewidth}L{0.22\linewidth}L{0.3\linewidth}L{0.24\linewidth}}
\toprule
位宽 $n$ & AND 门（部分积位） & 全加器（$n(n-1)$，逐行累加） & 关键路径 \\
\midrule
4 & 16 & 12 & 3 行加法链 \\
8 & 64 & 56 & 7 行加法链 \\
16 & 256 & 240 & 15 行加法链 \\
32 & 1024 & 992 & 31 行加法链 \\
\bottomrule
\end{longtable}

这就解释了为什么硬件乘法器比加法器贵得多：加法器是 $O(n)$ 个全加器排成一条链，乘法器是 $O(n^2)$ 个门铺成一片阵列，面积、延迟、功耗全面吃亏。所以小型控制器常用“每周期加一行”的移位加法方案，用时间换面积；只有乘法频繁出现的处理器才值得放一块完整阵列。现代高性能处理器进一步用 Booth 编码把部分积行数减半、用压缩树（如 Wallace 树）把逐行串行累加改成树形并行压缩，但“部分积矩阵加求和网络”这个基本结构没有改变。

\topic{Booth 编码用“一减一加”代替逐位累加}

逐位移位加法的累加次数等于乘数中 1 的个数，最坏情况下 n 位乘法要做 n 次。Booth 编码的观察是：一串连续的 1 可以整体改写。例如乘数中出现 $011110$，按位权展开是 $2^4+2^3+2^2+2^1=30$；但它也等于 $2^5-2^1=32-2$，也就是 $100000-000010$。原本 4 次加法的区段，改写成“高端加一次、低端减一次”就够了。二进制减法同样由加法器完成（取补相加），所以这种改写不引入新硬件，只减少累加次数。

两位 Booth 编码把这个想法系统化：从最低位开始，每次看乘数的当前位 $b_i$ 和它右边的相邻位 $b_{i-1}$（最低位右边补一个虚拟的 0），两位组合决定本轮对被乘数 $M$ 做什么：

\begin{longtable}{L{0.18\linewidth}L{0.18\linewidth}L{0.54\linewidth}}
\toprule
$b_i\,b_{i-1}$ & 本轮动作 & 含义 \\
\midrule
0 0 & $+0$ & 一串 0 的中间，无事可做 \\
0 1 & $+M$ & 一串 1 的开始（低端），把被乘数加上 \\
1 0 & $-M$ & 一串 1 的结束（高端），把被乘数减掉 \\
1 1 & $+0$ & 一串 1 的中间，已经折算过，跳过 \\
\bottomrule
\end{longtable}

直观记忆：从右往左扫乘数，0→1 是“1 串的入口”，做一次加；1→0 是“1 串的出口”，做一次减；00 和 11 都在串的内部，不累加。每轮做完后乘数右移一位，把下一对 bit 送上来。

以乘数 $0110$ 为例把滑动窗口画出来：右侧补一个虚拟 0，窗口从最低位起每次左移一位，四轮考察各做一次判断。

\begin{pdfdiagram}
  % 乘数 bit 行（右侧补虚拟 0）
  \draw[BookTeal, line width=0.9pt] (-0.55,2.6) rectangle (0.55,3.4);
  \node at (0,3.0) {0};
  \draw[BookTeal, line width=0.9pt] (0.65,2.6) rectangle (1.75,3.4);
  \node at (1.2,3.0) {1};
  \draw[BookTeal, line width=0.9pt] (1.85,2.6) rectangle (2.95,3.4);
  \node at (2.4,3.0) {1};
  \draw[BookTeal, line width=0.9pt] (3.05,2.6) rectangle (4.15,3.4);
  \node at (3.6,3.0) {0};
  \draw[BookTeal, line width=0.9pt] (4.25,2.6) rectangle (5.35,3.4);
  \node at (4.8,3.0) {0};
  \node[hw label] at (0,3.75) {$b_3$};
  \node[hw label] at (1.2,3.75) {$b_2$};
  \node[hw label] at (2.4,3.75) {$b_1$};
  \node[hw label] at (3.6,3.75) {$b_0$};
  \node[hw label] at (4.8,3.75) {补 0};
  % 扫描方向
  \draw[hw return, dashed] (5.3,4.2) -- (-0.5,4.2);
  \node[hw label, above] at (2.4,4.2) {窗口从右往左，每次左移一位};
  % 四轮滑动窗口
  \draw[BookAccent, line width=0.9pt, rounded corners=2pt] (3.05,1.95) rectangle (5.35,2.45);
  \node at (4.2,2.2) {$+0$};
  \node[hw label, anchor=east] at (-0.8,2.2) {轮 1};
  \node[hw label, anchor=west] at (5.6,2.2) {00：串外，无事可做};
  \draw[BookAccent, line width=0.9pt, rounded corners=2pt] (1.85,1.35) rectangle (4.15,1.85);
  \node at (3.0,1.6) {$-M$};
  \node[hw label, anchor=east] at (-0.8,1.6) {轮 2};
  \node[hw label, anchor=west] at (5.6,1.6) {10：1 串出口，减一次};
  \draw[BookAccent, line width=0.9pt, rounded corners=2pt] (0.65,0.75) rectangle (2.95,1.25);
  \node at (1.8,1.0) {$+0$};
  \node[hw label, anchor=east] at (-0.8,1.0) {轮 3};
  \node[hw label, anchor=west] at (5.6,1.0) {11：串中间，已折算};
  \draw[BookAccent, line width=0.9pt, rounded corners=2pt] (-0.55,0.15) rectangle (1.75,0.65);
  \node at (0.6,0.4) {$+M$};
  \node[hw label, anchor=east] at (-0.8,0.4) {轮 4};
  \node[hw label, anchor=west] at (5.6,0.4) {01：1 串入口，加一次};
  % 合计
  \node[hw label] at (2.4,-0.45) {合计：第 1 位权处减、第 3 位权处加，$8M-2M=6M$};
\end{pdfdiagram}
\diagramnote{两位 Booth 编码的滑动窗口（乘数 $0110$，右侧补虚拟 0）。窗口每次只看相邻两位 $b_i\,b_{i-1}$：0→1 是 1 串的入口，对应 $+M$；1→0 是出口，对应 $-M$；00 与 11 都在串的内部或串外，不累加。四轮中只有轮 2、轮 4 有动作，合计 $+M\cdot2^3-M\cdot2^1=8M-2M=6M$，与乘数值 6 一致——一连串 1 被折算成串尾之后的一次减和串首的一次加。}

以 4-bit 例子 $0011\times0110$（$3\times6=18$）走一遍。被乘数 $M=0011$，乘数 $B=0110$，右侧补 0 后逐对考察：

\begin{longtable}{L{0.12\linewidth}L{0.2\linewidth}L{0.16\linewidth}L{0.42\linewidth}}
\toprule
轮次 & 考察位 $b_i\,b_{i-1}$ & 动作 & 对累加结果的贡献 \\
\midrule
1 & 0 0 & $+0$ & 0 \\
2 & 1 0 & $-M$ & $-0011\times2^1=-6$ \\
3 & 1 1 & $+0$ & 0（1 串中间） \\
4 & 0 1 & $+M$ & $+0011\times2^3=+24$ \\
\bottomrule
\end{longtable}

合计 $-6+24=18$，结果正确。普通逐位法对乘数 $0110$ 要做 2 次加法，这个例子里 Booth 也是一加一减，看似没有节省。差别在 1 串变长时才显现：乘数若是 $011110$（4 个连续 1），逐位法要 4 次加法，Booth 仍只要“低端一加、高端一减”共 2 次动作；乘数若是全 1 的 $1111$，逐位法要 4 次加法。这里要分清两种读法。按补码读，$1111$ 就是 $-1$，Booth 只在最低位遇到一次 1→0（实际是考察位 (1,0)），做一次减就结束，$-M$ 正是 $M\times(-1)$。若按无符号把它读成 15，则要在左侧再补一个 0、按 $01111$ 扫描：最低位的一次减之外，补入的 0 与串首的 1 又构成一次 0→1，串首之后还有一次加，合计 $16M-M=15M$，仍只有两次动作。

Booth 编码的意义不在于硬件更精巧，而在于把“乘法要做几次累加”从“取决于乘数里有几个 1”变成“取决于乘数里有几段 1”。它还附带一个好处：这套规则天然兼容补码有符号乘数，不需要像原码乘法那样单独处理符号位。现代乘法器普遍采用按两位、三位分组的改进 Booth（radix-4、radix-8），把部分积行数减到 $n/2$ 或 $n/3$ 行，再配合压缩树求和——这就是乘法器能在一个流水段里算完的起点。

\topic{恢复余数除法：一次一位地试出商}

除法是乘法的逆过程，但硬件上它比乘法更难安排：乘法每一步做什么由乘数位直接决定，除法每一步要先“试一试”才知道。恢复余数除法（restoring division）是最直白的方案，和笔算长除法同构：每轮把部分余数左移一位，试减除数；够减，这一位商就是 1，余数保留差值；不够减，这一位商是 0，把余数恢复成试减前的值，进入下一轮。

以 4-bit 无符号除法 $13\div3$ 为例：被除数 $1101$，除数 $0011$。用一个 8-bit 寄存器 $R$ 存放部分余数，高 4 位初始为 0，低 4 位初始为被除数，即 $R=0000\,1101$；除数始终与高 4 位对齐比较。每轮三步：$R$ 整体左移一位；高 4 位试减 $0011$；差不为负则保留差并把最低位置 1，否则恢复高 4 位、最低位置 0。

\begin{longtable}{L{0.12\linewidth}L{0.22\linewidth}L{0.28\linewidth}L{0.1\linewidth}L{0.18\linewidth}}
\toprule
轮次 & 左移后的 $R$ & 高 4 位试减 $0011$ & 商位 & 本轮结束的 $R$ \\
\midrule
初始 & --- & --- & --- & 0000 1101 \\
1 & 0001 1010 & $0001-0011$，不够减，恢复 & 0 & 0001 1010 \\
2 & 0011 0100 & $0011-0011=0000$，够减 & 1 & 0000 0101 \\
3 & 0000 1010 & $0000-0011$，不够减，恢复 & 0 & 0000 1010 \\
4 & 0001 0100 & $0001-0011$，不够减，恢复 & 0 & 0001 0100 \\
\bottomrule
\end{longtable}

四轮结束后，低 4 位 $0100$ 是商，高 4 位 $0001$ 是余数：$13=3\times4+1$，与笔算一致。注意第 2 轮左移后，被除数的高位随左移进入高 4 位，恰好第一次拼出够减的 $0011$——这和笔算长除法里“逐位落数，直到够减”是同一个过程。

把每轮的三个动作画成数据通路，硬件只剩三块：一个 (R,Q) 拼接移位寄存器、一个试减器、一个按符号选择的写回 mux。

\begin{pdfdiagram}
  % 顶部：拼接移位寄存器与除数寄存器
  \node[hw state, minimum width=3.6cm] (rq) at (-2.2,3.2) {(R,Q) 拼接\\移位寄存器};
  \node[hw state, minimum width=1.8cm] (dv) at (2.6,3.2) {除数 D};
  % 中部：试减器与符号判负
  \node[hw compute, minimum width=3.2cm] (sub) at (0.2,0.6) {试减器 R 高 $-$ D};
  \node[hw control, minimum width=3.0cm] (sgn) at (4.6,0.6) {符号判负\\差 $<0$？};
  % 底部：条件提交 mux
  \node[hw compute, minimum width=3.4cm] (mux) at (0.2,-1.6) {条件提交 mux\\差值 / 原值};
  % 路径①：左移自环
  \draw[hw bus] (-0.9,3.72) -- (-0.9,4.5) -- (-3.5,4.5) -- (-3.5,3.72);
  \node[hw label] at (-2.2,4.85) {① 左移：(R,Q) 左移 1 位};
  % 路径②：试减的两路输入
  \draw[hw bus] (-0.9,2.68) -- (-0.9,1.7) -- (-0.45,1.7) -- (-0.45,1.0);
  \node[hw label, anchor=west] at (-0.75,2.2) {R 高 4 位};
  \draw[hw bus] (2.6,2.8) -- (2.6,1.7) -- (0.85,1.7) -- (0.85,1.0);
  \node[hw label, anchor=east] at (2.45,2.3) {除数 D};
  \node[hw label, anchor=east] at (-1.6,0.6) {② 试减};
  % 差值送符号判负，判负结果（虚线）选通 mux
  \draw[hw bus] (1.8,0.6) -- (3.0,0.6);
  \node[hw label, above] at (2.4,0.6) {差值};
  \draw[hw return, dashed] (4.6,0.08) -- (4.6,-1.6) -- (1.9,-1.6);
  \node[hw label, above] at (3.1,-1.6) {③ 判负选通};
  % 试减差值下行进 mux
  \draw[hw bus] (0.2,0.2) -- (0.2,-1.08);
  \node[hw label, anchor=west] at (0.35,-0.45) {差值};
  % 恢复路径：原值绕左侧进 mux
  \draw[hw bus] (-4.0,3.2) -- (-4.6,3.2) -- (-4.6,-0.7) -- (-1.0,-0.7) -- (-1.0,-1.08);
  \node[hw label, rotate=90] at (-5.0,1.2) {原值（恢复）};
  % 路径③：条件提交回写 R 高 4 位
  \draw[hw bus] (0.2,-2.12) -- (0.2,-2.9) -- (-3.4,-2.9) -- (-3.4,2.68);
  \node[hw label, below] at (-1.6,-2.9) {③ 条件提交：回写 R 高 4 位};
\end{pdfdiagram}
\diagramnote{恢复余数除法的数据通路：每轮三条路径。① (R,Q) 整体左移一位，空出的 Q 低位留给本轮商位；② R 高 4 位与除数试减；③ 按差的符号条件提交——差非负则写回差值、商位记 1，差为负则沿左侧恢复路径写回原值、商位记 0。“不够减要恢复”在图上就是 mux 多出来的那路原值输入：恢复不需要再做一次加法，只是把试减前的旧值选回来。}

为什么除法非要 n 轮不可：商有 n 位，每一位都要独立判断一次“够不够减”，而这个判断的输入是上一轮算完的余数，无法并行预知。所以除法本质上是串行的，n 位商就要 n 轮，每轮一次减法（同时就是比较）加一次条件提交。恢复余数法的代价是那一次“白做的减法”：不够减时要恢复，硬件上体现为每轮多一个二选一（减的结果与原值），或者多一步把原值加回来。

不恢复余数法（non-restoring）发现“本轮恢复、下轮再减”可以合并成“下轮直接加”，省掉恢复动作；SRT 除法（得名于 Sweeney、Robertson 和 Tocher）用商数字集合 $\{-1,0,+1\}$ 加查找表，一轮产生多位商。但无论哪种方案，“每轮一次比较、一次条件更新余数”的串行骨架不变——这就是现代处理器里整数除法仍要十几到几十个周期、而乘法只要几个周期的根本原因。

\begin{chipcase}
74181 是经典 4-bit ALU 芯片，常被用来观察算术逻辑单元如何把多种函数、进位输入、进位输出和模式控制组合到一个器件中。它不是现代 CPU 的内部实现模板，却非常适合作为教学标本：控制引脚选择函数，数据引脚提供操作数，输出引脚给出结果和进位信息。把 74181 与 8086、80386 对照，可以看到同一问题的放大版本：位宽增加、控制复杂度上升、标志位更多地参与分支、地址和异常判断，ALU 不再只是一个孤立芯片，而是数据通路和控制器中的高频核心。
\end{chipcase}

用 74181 做实验时，可以把每一类引脚对应到本章讲过的控制对象：A/B 是操作数输入，F 是结果输出，S 选择具体函数，M 选择逻辑模式或算术模式，\code{Cn} 是低位进位输入，\code{Cn+4} 是高位进位输出，若要扩展到更宽位宽，还要用跨片进位或专门的超前进位芯片处理片间等待。

\begin{longtable}{L{0.18\linewidth}L{0.36\linewidth}L{0.36\linewidth}}
\toprule
74181 观察点 & 接线含义 & 对应本章概念 \\
\midrule
A/B 输入 & 两个 4-bit 操作数进入芯片 & ALU 操作数总线 \\
S 控制 & 选择具体逻辑或算术函数 & \code{alu_op} 的一部分 \\
M 控制 & 区分逻辑类函数和算术类函数 & ALU 模式选择 \\
\code{Cn} & 低位进位输入，可接前一级进位或控制信号 & 加法、减法、多字长运算入口 \\
\code{Cn+4} & 4-bit 组的高位进位输出 & Carry 标志或级联输入 \\
F 输出 & 4-bit 结果输出 & 写回候选结果 \\
\bottomrule
\end{longtable}

若用两片 74181 组成 8-bit ALU，低 4 位芯片处理 bit0 到 bit3，高 4 位芯片处理 bit4 到 bit7。最低位 \code{Cn} 由本次动作决定：普通加法为 0，带进位加法接 Carry，减法按芯片函数和控制约定接入相应初始进位。低片的 \code{Cn+4} 送到高片 \code{Cn}，即可得到串行跨片进位。若继续扩展到 16-bit，这种片间进位会增加等待；因此经典板级设计常配合 74182 一类超前进位发生器，减少多片 74181 之间逐片等待的时间。更宽的数据通路不是简单复制芯片，还要重新处理进位、控制和时序。

若用分立芯片搭一个 4-bit ALU 实验台，可以按以下顺序推进。第一步只搭 74HC283 加法；第二步加入 74HC86 实现加/减复用；第三步用 74HC08、74HC32、74HC86 分别产生 AND、OR、XOR 候选结果；第四步用 74HC157 选择加法结果或逻辑结果；第五步用 LED 显示 Zero、Carry、Negative 和 Overflow。每一步都要确保前一步的结果仍可观察、可复现，避免在复杂电路里同时排查过多变量。

\begin{longtable}{L{0.2\linewidth}L{0.38\linewidth}L{0.32\linewidth}}
\toprule
实验阶段 & 新增器件或连线 & 预期现象 \\
\midrule
4-bit 加法 & 74HC283、拨码开关、LED & \code{0011+0101=1000}，Carry 正确 \\
加/减复用 & 74HC86 处理 B，\code{sub} 接 \code{Cin} & \code{sub=1} 时 \code{A-B} 正确 \\
逻辑候选 & AND/OR/XOR 门阵列 & 候选结果与真值表一致 \\
结果选择 & 74HC157 多路选择器 & \code{op} 改变时只有被选结果进入 LED \\
标志输出 & OR 合成 Zero，最高位和进位生成标志 & 结果为 0 时 Zero=1，溢出案例可复现 \\
\bottomrule
\end{longtable}
